\chapter{Compact Objects \& Relativistic Astrophysics}

\section*{Theory problems}

\probhead{10.1}{}{2007}{Chiang Mai, Thailand}{TH 1.14}{2}
% source id: 2007_TH_1.14   tag: Schwarzschild radius of the Sun
If the Sun were to collapse gravitationally to form a non-rotating black hole,
what would be its event horizon (its Schwarzschild radius)?

\probhead{10.2}{Gravitational Lensing}{2007}{Chiang Mai, Thailand}{TH Q4}{10}
% source id: 2007_TH_Q4   tag: Einstein ring, two lensed images
\ioaafig{2007_TH_Q4-f1.pdf}

The deflection of light by a gravitational field was first predicted by
Einstein in 1912 a few years before the publication of the General Relativity
in 1916. A massive object that causes a light deflection behaves like a
classical lens. This prediction was confirmed by Sir Arthur Stanley Eddington
in 1919.

\ioaafig{2007_TH_Q4-f2.pdf}

Consider a spherically symmetric lens, with a mass $M$ with an impact
parameter $\xi$ from the centre. The deflection equation in this case is given
by:
\[ \phi = \frac{4GM}{\xi c^2}, \]
a very small angle. In figure 1, the massive object which behaves like a lens
is at L. Light rays emitted from the source S being deflected by the lens are
observed by observer O as images $\mathrm{S}_1$ and $\mathrm{S}_2$. Here,
$\phi, \beta$, and $\theta$ are very very small angles.

\begin{description}
\item[a)] For a special case in which the source is perfectly aligned with the
  lens such that $\beta = 0$, show that a ring-like image will occur with the
  angular radius, called Einstein radius $\theta_E$, given by:
  \[ \theta_E = \sqrt{\left(\frac{4GM}{c^2}\right)
     \left(\frac{D_S - D_L}{D_L D_S}\right)} \]
  \hfill[2\,pt]
\item[b)] The distance (from Earth) to a source star is about 50\,kpc. A
  solar-mass lens is about 10\,kpc from the star. Calculate the angular radius
  of the Einstein ring formed by this solar-mass lens with the perfect
  alignment. \hfill[1\,pt]
\item[c)] What is the resolution of the Hubble space telescope with 2.4\,m
  diameter mirror? Could the Hubble telescope resolve the Einstein ring in
  b)? \hfill[2\,pt]
\item[d)] In figure 1, for an isolated point source S, there will be two
  images ($\mathrm{S}_1$ and $\mathrm{S}_2$) formed by the gravitational
  lens. Find the positions ($\theta_1$ and $\theta_2$) of the two images.
  Answer in terms of $\beta$ and $\theta_E$. \hfill[2\,pt]
\item[e)] Find the ratio $\frac{\theta_{1,2}}{\beta}$ ($\frac{\theta_1}{\beta}$
  or $\frac{\theta_2}{\beta}$) in terms of $\eta$. Here $\theta_{1,2}$
  represents each of the image positions in d) and $\eta$ stands for the ratio
  $\frac{\beta}{\theta_E}$. \hfill[2\,pt]
\item[f)] Find also the values of magnifications
  $\frac{\Delta\theta}{\Delta\beta}$ in terms of $\eta$ for
  $\theta = \theta_{1,2}$ ($\theta = \theta_1$ or $\theta = \theta_2$), when
  $\Delta\beta \ll \beta$, and $\Delta\theta \ll \theta$. \hfill[1\,pt]
\end{description}

\probhead{10.3}{}{2009}{Tehran, Iran}{TH P1}{10}
% source id: 2009_TH_P1   tag: SMBH mean density inside Schwarzschild radius
Calculate the mean mass density for a super massive black hole with total mass
of $1 \times 10^8\,M_\odot$ inside the Schwarzschild radius.

\probhead{10.4}{}{2009}{Tehran, Iran}{TH P10}{10}
% source id: 2009_TH_P10   tag: magnetic flux conservation, neutron star
A main sequence star with the radius and mass of $R = 4R_\odot$,
$M = 6M_\odot$ has an average magnetic field of $1 \times 10^{-4}$\,T.
Calculate the strength of the magnetic field of the star when it evolves to a
neutron star with the radius of 20\,km.

\probhead{10.5}{}{2010}{Beijing, China}{TH S2}{10}
% source id: 2010_TH_S2   tag: radius where escape velocity exceeds light
If the escape velocity from a solar mass object's surface exceeds the speed of
light, what would be its radius?

\probhead{10.6}{}{2012}{Rio de Janeiro, Brazil}{TH S8}{}
% source id: 2012_TH_S8   tag: pulsar beam detection probability
A pulsar located 1000\,pc far from Earth, 10\,000 times more luminous than our
Sun, emits radiation only from its two opposite poles, creating an homogeneous
emission beam shaped as double cone with opening angle $\alpha = 4^\circ$.
Assuming the angle between the rotation axis and the emission axis is
$30^\circ$, and assuming a random orientation of the pulsar beams in relation
to an observer on Earth, what is the probability of detecting the pulses? In
case we can see it, what is the apparent bolometric magnitude of the pulsar?
\ioaafig{2012_TH_S8-f1.pdf}

\probhead{10.7}{Cosmic radiation}{2014}{Suceava, Romania}{TH P3}{}
% source id: 2014_TH_P3   tag: pion decay relativistic kinematics
During studies concerning cosmic radiation, a neutral unstable particle -- the
$\pi^0$ meson was identified. The rest-mass of meson $\pi^0$ is much larger
than the rest-mass of the electron. The studies reveal that during its flight,
the meson $\pi^0$ disintegrates into 2 photons.

Find an expression for the initial velocity of the meson $\pi^0$, if after its
disintegration, one of the photons has the maximum possible energy
$E_{\mathrm{max}}$ and, consequently, the other photon has the minimum
possible energy $E_{\mathrm{min}}$. You may use as known $c$ -- the speed of
light.

\probhead{10.8}{}{2015}{Magelang, Indonesia}{TH S12}{}
% source id: 2015_TH_S12   tag: relativistic proton Sun-Earth travel time
A proton with a kinetic energy of $1\,\mathrm{GeV}$ propagates out from the
surface of the Sun towards the Earth. Neglecting the magnetic field of the
Sun, calculate the travel time of the proton as seen from the Earth.

\probhead{10.9}{Gravitational Lensing Telescope}{2016}{Bhubaneswar, India}{TH T10}{20}
% source id: 2016_TH_T10   tag: solar gravitational lens focus
Einstein's General Theory of Relativity predicts bending of light around
massive bodies. For simplicity, we assume that the bending of light happens at
a single point for each light ray, as shown in the figure. The angle of
bending, $\theta_{\mathrm{b}}$, is given by
\[ \theta_{\mathrm{b}} = \frac{2R_{\mathrm{sch}}}{r} \]
where $R_{\mathrm{sch}}$ is the Schwarzschild radius associated with that
gravitational body. We call $r$, the distance of the incoming light ray from
the parallel $x$-axis passing through the centre of the body, as the ``impact
parameter''.

\ioaafig{2016_TH_T10-f1.pdf}

A massive body thus behaves somewhat like a focusing lens. The light rays
coming from infinite distance beyond a massive body, and having the same
impact parameter $r$, converge at a point along the axis, at a distance $f_r$
from the centre of the massive body. An observer at that point will benefit
from huge amplification due to this gravitational focusing. The massive body
in this case is being used as a Gravitational Lensing Telescope for
amplification of distant signals.

\begin{description}
\item[(T10.1)] Consider the possibility of our Sun as a gravitational lensing
  telescope. Calculate the shortest distance, $f_{\mathrm{min}}$, from the
  centre of the Sun (in A.U.) at which the light rays can get focused.
  \hfill[6\,pt]
\item[(T10.2)] Consider a small circular detector of radius $a$, kept at a
  distance $f_{\mathrm{min}}$ centered on the $x$-axis and perpendicular to
  it. Note that only the light rays which pass within a certain annulus (ring)
  of width $h$ (where $h \ll R_\odot$) around the Sun would encounter the
  detector. The amplification factor at the detector is defined as the ratio
  of the intensity of the light incident on the detector in the presence of
  the Sun and the intensity in the absence of the Sun.

  Express the amplification factor, $A_{\mathrm{m}}$, at the detector in terms
  of $R_\odot$ and $a$. \hfill[8\,pt]
\item[(T10.3)] Consider a spherical mass distribution, such as dark matter in
  a galaxy cluster, through which light rays can pass while undergoing
  gravitational bending. Assume for simplicity that for the gravitational
  bending with impact parameter, $r$, only the mass $M(r)$ enclosed inside the
  radius $r$ is relevant.

  What should be the mass distribution, $M(r)$, such that the gravitational
  lens behaves like an ideal optical convex lens? \hfill[6\,pt]
\end{description}

\probhead{10.10}{Deflection of radio photons in the gravitational field of solar system bodies}{2019}{Keszthely, Hungary}{TH 2}{10}
% source id: 2019_TH_2   tag: GR light deflection by Jupiter/Moon
A.~Eddington and F.~Dyson from Principe, and C.~Davidson and A.~Crommelin from
Sobral, Brazil measured the deflection of light coming from stars apparently
very close to the Sun during the total solar eclipse in 1919. The deflection
was found to agree with the theoretically predicted value of $1.75''$.

A light ray (or photon) which passes the Sun at a distance $d$ is deflected by
an angle
\[ \Delta\theta \propto \frac{4GM_\odot}{d c^2} \]

The present-day accuracy of the VLBI (Very Long Baseline Interferometry)
technique in the radio wavelength range is 0.1\,mas (milliarcsecond). Is it
possible to detect a deflection of radio photons from a quasar by (a) Jupiter,
(b) the Moon? Estimate the angle of deflection in both cases and mark ``YES''
or ``NO'' on the answer sheet.

\probhead{10.11}{The supermassive black hole in the centre of Milky Way Galaxy and M87}{2019}{Keszthely, Hungary}{TH 3}{10}
% source id: 2019_TH_3   tag: EHT resolving BH shadow
The first image of a black hole was constructed recently by the international
team of the Event Horizon Telescope (EHT). The imaged area surrounds the
supermassive black hole in the centre of the galaxy M87. The observations
resulting in the final image were carried out at a wavelength
$\lambda = 1.3$\,mm where the interstellar extinction is not prohibitively
large.

\begin{description}
\item[a)] How large an instrument would be needed to resolve the shadow (in
  effect the photon capture radius, which is three times the size of event
  horizon) of a supermassive black hole in the centre of a galaxy? Express the
  result as a function of the distance $d$ and the mass $M$ of the black hole.
  \hfill[6\,p]
\item[b)] Give the size of the instrument in units of Earth radius for
  \begin{description}
  \item[i.] the supermassive black hole in the centre of M87
    ($d_{\mathrm{BH\text{-}M87}} = 5.5 \times 10^7$\,ly,
    $M_{\mathrm{BH\text{-}M87}} = 6.5 \times 10^9\,\mathrm{M}_\odot$),
    \hfill[1\,p]
  \item[ii.] and Sgr~A*, the supermassive black hole of our own galaxy, Milky
    Way ($d_{\mathrm{Sgr\,A^*}} = 8.3$\,kpc,
    $M_{\mathrm{Sgr\,A^*}} = 3.6 \times 10^6\,\mathrm{M}_\odot$). \hfill[1\,p]
  \end{description}
\item[c)] What type of technology is needed for the development of such an
  instrument? Mark the letter of your answer with $\times$ on the answer
  sheet. \hfill[2\,p]
  \begin{description}
  \item[(A)] Gravitational lensing by dark matter
  \item[(B)] Interferometry with an array of radio telescopes
  \item[(C)] Photon deceleration in a dense environment
  \item[(D)] Reducing the effect of incoming wavefront distortions
  \item[(E)] Neutrino focusing with strong electromagnetic fields
  \end{description}
\end{description}

\probhead{10.12}{Planck's units}{2022}{Kutaisi, Georgia}{TH 1}{10}
% source id: 2022_TH_1   tag: Planck-scale dimensional analysis
In Max Planck's system of natural units of measurements, commonly used in
cosmology, all units are expressed in terms of 4 fundamental constants: speed
of light ($c$), universal gravitational constant ($G$), reduced Planck
constant ($\hbar = \frac{h}{2\pi}$) and Boltzmann constant ($k_B$). Also,
$(4\pi\epsilon_0)$ is included in this list of constants.

Using dimensional analysis, find expressions for
\begin{itemize}
\item Planck length ($L_p$),
\item Planck time ($t_p$),
\item Planck mass ($m_p$),
\item Planck temperature ($T_p$) and
\item the Planck charge ($q_p$).
\end{itemize}

\probhead{10.13}{Accretion}{2022}{Kutaisi, Georgia}{TH 10}{20}
% source id: 2022_TH_10   tag: Eddington limit, accretion luminosity
Consider a compact object (such as black hole, white dwarf or neutron star)
with a spherically symmetric accretion of gas, assume that accreting gas is
hydrogen. As particles fall into the object they heat up and radiate, thus
creating radiation pressure acting on rest of the accreting material. This
force is given by,
\[ F_L = \sigma_e \frac{I}{c} \]
where,
\[ \sigma_e = \frac{8\pi}{3}
   \left(\frac{e^2}{4\pi\epsilon_0 m_e c^2}\right)^2 \]
is the Thompson cross section for electrons, $c$ is the speed of light and $I$
intensity of the light. Although $F_L$ is calculated for electrons, it
effectively acts on the whole atom.

\begin{description}
\item[(a)] If the central compact object has a mass $M$, find an expression
  for the Eddington limit ($L_E$), which is the maximum possible luminosity
  for the accretion sphere. \hfill[3\,pt]
\item[(b)] For a particular compact object, the luminosity due to material
  accretion is the same as the Solar luminosity $L_\odot$. What is its minimum
  possible mass of this object to achieve this luminosity (in units of
  $M_\odot$)? \hfill[4\,pt]
\end{description}

Assume that the atoms in the accretion sphere originate far away from the
compact object. When these atoms fall into the compact object, their
gravitational energy is transformed into radiation.

\begin{description}
\item[(c)] Derive an expression for accretion luminosity ($L_{\mathrm{acc}}$)
  in terms of the compact object's mass ($M$), mass accretion rate
  ($\dot{M} = \frac{\Delta M}{\Delta t}$), and the compact object's radius
  ($R$). \hfill[2\,pt]
\item[(d)] Show that the maximum possible accretion rate in steady state is
  not directly dependent on the mass of the compact object. \hfill[2\,pt]
\item[(e)] For an object with $R = 12 \times 10^9$\,m, calculate the maximum
  possible accretion rate $\dot{M}$ (in units of $M_\odot$ per year).
  \hfill[2\,pt]
\end{description}

In reality, the accretion geometry is disk shaped, where most particles trace
an almost circular orbit around the compact object. Consider a binary system
consisting of a compact object of mass $M_1$ and a hydrogen burning star of
mass $M_2$ at a distance $a$ from each other. The gas from the hydrogen
burning star is accreted by the compact object and due to this mass transfer
the period of the binary system changes.

\begin{description}
\item[(f)] Find the condition on the two masses such that the separation
  between the stars is increasing. Ignore the rotation of the stars.
  \hfill[7\,pt]
\end{description}

\probhead{10.14}{Relativistic Beaming}{2022}{Kutaisi, Georgia}{TH 13}{50}
% source id: 2022_TH_13   tag: relativistic Doppler, aberration, M87 jet
Consider an isotropic light source of frequency $f_R$ in a frame which is
fixed to the source (i.e.\ rest frame). In this rest frame, consider a light
ray emitted from the source that makes an angle $\theta_R$ with the $X$-axis.
The light source is moving along positive $X$ direction with (relativistic)
speed $v$ as measured in the lab frame.

\begin{description}
\item[(a)] Find an expression for the frequency $f_L$ of this ray in the lab
  frame, and the cosine of the angle that this ray makes with the $X$-axis in
  the lab frame. \hfill[11\,pt]\\
  \textit{Hint:} In relativistic mechanics, energy $E$ and momentum $p$ of a
  particle between rest and lab frame are related in the following way:
  \begin{align*}
    \frac{E_L}{c} &= \gamma\left(\frac{E_R}{c} + p_{x_R}\frac{v}{c}\right)\\
    p_{x_L} &= \gamma\left(p_{x_R} + \frac{E_R v}{c^2}\right)\\
    p_{y_L} &= p_{y_R}\\
    p_{z_L} &= p_{z_R}
  \end{align*}
  where:
  \[ \gamma = \frac{1}{\sqrt{1 - \frac{v^2}{c^2}}} \]
\item[(b)] For the following cases:
  \begin{description}
  \item[i)] $\theta_R = 0^\circ$
  \item[ii)] $\theta_R = \cos^{-1}(-v/c)$
  \item[iii)] $\theta_R = 90^\circ$
  \item[iv)] $\theta_L = 180^\circ$
  \end{description}
  draw direction vectors of the beam in $XY$ plane of the rest frame as well
  as separately in $X'Y'$ plane of the lab frame. \hfill[4\,pt]
\end{description}

In accretion disks around black holes, the charged particles are orbiting at
relativistic speeds and in their rest frames may be considered as isotropic
point sources of light. Consider such a particle K in a circular orbit of
radius $r$ and angular speed $\omega$ around a central object located at $O$
(see figure).

\ioaafig{2022_TH_13-f1.pdf}

Let us assume that our lab frame is fixed to an observer located at a point
$P$ on the $OY$ axis, which is stationary with respect to $O$.
$OP = R \gg r$. Let $t_{L0} = t_{R0} = 0$ correspond to the moment when K is
seen crossing the $OX$ axis. As K is moving with relativistic speed, the
duration $\Delta t_R$ measured by an observer in the rest frame of the source
K is related to the duration measured in the lab frame $\Delta t_L$ at $P$ by
the expression $\Delta t_L = \gamma\Delta t_R$.

\begin{description}
\item[(c)] Derive an expression for $f_L$ as a function of $t_L$
  ($t_L > R/c$)? \hfill[7\,pt]
\end{description}

Let us consider a fraction of the light from the source that is emitted in an
infinitesimal solid angle
$\Delta\Omega_R = -\Delta(\cos\theta_R) \cdot \Delta\phi$ in the direction
making an angle $\theta_R$ with respect to the $X$ axis in the rest frame, as
it is shown on the figure below.

\ioaafig{2022_TH_13-f2.pdf}

\begin{description}
\item[(d)] Show that, as measured in the lab frame
  \[ \Delta\Omega_L = \frac{\Delta\Omega_R}
     {\gamma^2\left(1 + \frac{v}{c}\cos\theta_R\right)^2} \]
  \hfill[10\,pt]
\item[(e)] If the intrinsic luminosity of the light source is $L$, what is the
  energy flux $F_L$ observed by the observer at point $P$ at the moment $t_L$
  ($t_L > R/c$)? \hfill[15\,pt]\\
  \textit{Hint:} In the rest frame of the source, you may assume $N_R$ number
  of photons are directed within the solid angle $\Delta\Omega_R$ during the
  time interval $\Delta t_R$.
\item[(f)] Charged particles in the relativistic beam shot from the
  supermassive black hole at the centre of the galaxy M87 have speeds up to
  $0.95c$. What would be the maximum and minimum amplification factor for the
  energy flux for a relativistic beam from M87? \hfill[3\,pt]
\end{description}

\probhead{10.15}{White Dwarf}{2024}{Rio de Janeiro, Brazil}{TH T4}{10}
% source id: 2024_TH_T4   tag: degenerate electron gas mass-radius
The structure of a white dwarf is sustained against gravitational collapse by
the pressure of degenerate electrons, a phenomenon explained by quantum
physics and related to the Pauli Exclusion Principle for electrons. The
equation of state of a gas made of non-relativistic degenerate electrons is
the following:
\[ P = \left(\frac{3}{8\pi}\right)^{2/3}\frac{h^2}{5m_e}\,n_e^{5/3}, \]
where $n_e$ is the number of electrons per unit volume, which can be expressed
in terms of the mass density $\rho$ using the dimensionless factor $\mu_e$,
the number of nucleons (protons and neutrons) per unit electron. Also consider
that the central pressure can be described by this equation of state.

In the condition of hydrostatic equilibrium, the pressure and gravitational
forces balance each other at any distance $r$ from the centre of the star.
This condition can be expressed by:
\[ \frac{\mathrm{d}P}{\mathrm{d}r} = -\frac{GM(r)\rho(r)}{r^2}, \]
where $M(r)$ is the mass contained in the sphere of radius $r$, and $\rho(r)$
is the mass density of the star at a radius $r$.

Assume that $m_p = m_n$, the density of a white dwarf is roughly uniform, and
the following approximation is valid at the surface of the star:
\[ \left.\frac{\mathrm{d}P}{\mathrm{d}r}\right|_{r=R}
   \approx -\frac{P_c}{R}, \]
where $P_c$ is the pressure at the centre of the star, and $R$ the star
radius.

\begin{description}
\item[(a)] The relationship between the mass $M$ and the radius $R$ of a white
  dwarf can be written in the form:
  \[ R = a \cdot M^b \]
  Find the exponent $b$ and determine the coefficient $a$ in terms of physical
  constants and $\mu_e$. \hfill[6\,pt]
\item[(b)] Using the relationship found in the previous part, estimate the
  radius of a white dwarf made of fully ionised carbon
  ($^{12}_{6}\mathrm{C}$) with a mass of $M = 1.0\,M_\odot$. \hfill[4\,pt]
\end{description}

\probhead{10.16}{Physics of Accretion}{2024}{Rio de Janeiro, Brazil}{TH T9}{35}
% source id: 2024_TH_T9   tag: thin accretion disc luminosity
The accretion of matter onto compact objects, such as neutron stars and black
holes, is one of the most efficient ways to produce radiant energy in
astrophysical systems. Consider an element of gas of mass $\Delta m$ in a
stationary and geometrically thin disc of matter with a maximum radius of
$R_{\mathrm{max}}$ and minimum stable orbital radius of $R_{\mathrm{min}}$
(with $R_{\mathrm{min}}/R_{\mathrm{max}} \ll 1$), in rotation around a compact
object of mass $M$ and radius $R$.

\begin{description}
\item[(a)] Assuming that an element of gas in the disc follows an
  approximately Keplerian circular orbit, find an expression for the total
  mechanical energy per unit mass $\frac{\Delta E}{\Delta m}$ released by this
  gas from when it starts orbiting at a radius $R_{\mathrm{max}}$ until it
  reaches an orbital radius of $r \ll R_{\mathrm{max}}$. This process occurs
  very slowly, transforming kinetic energy into internal energy of the gas
  disc through viscous dissipation. \textbf{Note:} Ignore the gravitational
  interaction between particles within the accretion disc and give your final
  answer in terms of $G$, $M$, and $r$. \hfill[6\,pt]
\item[(b)] Considering that the disc receives mass at an average rate of
  $\dot{M}$, and assuming that all the mechanical energy lost is ultimately
  converted into radiation, find an expression for the total luminosity of the
  disc. \hfill[5\,pt]
\item[(c)] Consider now the ring composed of all mass elements with radii
  between $r$ and $r + \Delta r$. In this scenario, find an expression for the
  luminosity generated by the disc over its small width $\Delta r$ at this
  radius; that is, find the expression for
  $\frac{\Delta E}{\Delta t\,\Delta r}$. \hfill[8\,pt]
\item[(d)] Assuming that the gravitational energy released in this ring is
  locally emitted by the surface of the ring in the form of black-body
  radiation, find an expression for the surface temperature $T$ of the ring.
  \hfill[10\,pt]
\item[(e)] Consider that the central object is a stellar black hole with a
  mass of $3M_\odot$ and a rate of accretion of
  $\dot{M} = 10^{-9}\,M_\odot/\mathrm{year}$. Consider also that
  $R_{\mathrm{min}} = 3R_{\mathrm{sch}}$, where $R_{\mathrm{sch}}$ is the
  Schwarzschild radius of the black hole. Determine the luminosity of the disc
  and the peak wavelength of emission of its innermost part. Ignore
  gravitational redshift effects and assume that the emission from the
  innermost part of the ring dominates the total emission. \hfill[3\,pt]
\item[(f)] Now, considering another accretion system with
  $\dot{M} = 1\,M_\odot/\mathrm{year}$ and a peak emission wavelength of
  $\lambda = 6 \times 10^{-8}$\,m, estimate the mass of this black hole.
  \hfill[3\,pt]
\end{description}

\probhead{10.17}{Neutron Star Binary}{2025}{Mumbai, India}{TH T07}{20}
% source id: 2025_TH_T07   tag: wind-fed accretion, magnetospheric radius
In a binary system involving a compact star, where the companion star does not
overflow its Roche lobe, the primary source of accretion for the compact star
is the stellar wind from the companion star. This wind-fed accretion is
especially significant in systems that include an early-type star (such as an
O or B star, indicated henceforth as an OB star), alongside a compact object
like a neutron star (NS) in a close orbit.

Consider such a NS-OB star binary system where a neutron star of mass
$M_{\mathrm{NS}} = 2.0\,\mathrm{M}_\odot$ and radius
$R_{\mathrm{NS}} = 11$\,km is orbiting in a circular orbit of radius $a$
around the centre of the OB star with velocity
$v_{\mathrm{orb}} = 1.5 \times 10^5\,\mathrm{m\,s^{-1}}$ (see figure below).
Throughout this problem the mass loss from the OB star is assumed to be
spherically symmetric and its rate is
$\dot{M}_{\mathrm{OB}} = 1.0 \times 10^{-4}\,\mathrm{M}_\odot\,
\mathrm{yr}^{-1}$.

\ioaafig{2025_TH_T07-f1.pdf}

\begin{description}
\item[(T07.1)] The accretion radius, $R_{\mathrm{acc}}$, is defined as the
  maximum distance from the NS at which the stellar wind can be captured by
  the NS. If the stellar wind speed at the orbital distance of the NS is
  $v_{\mathrm{w}} = 3.0 \times 10^6\,\mathrm{m\,s^{-1}}$, find
  $R_{\mathrm{acc}}$ for the above system in km using standard escape velocity
  calculations. \hfill[3\,pt]
\item[(T07.2)] Assuming that all captured material is accreted by the NS,
  estimate the mass accretion rate, $\dot{M}_{\mathrm{acc}}$, from the stellar
  wind onto the NS in units of $\mathrm{M}_\odot\,\mathrm{yr}^{-1}$ if
  $a = 0.5$\,au. Neglect the effects of radiation pressure and finite cooling
  time of the accreting gas. \hfill[3\,pt]
\item[(T07.3)] Now consider the situation where the stellar wind speed at the
  orbital distance $a$ (near the NS) becomes comparable with orbital speed of
  the NS. The mass accretion rate from the stellar wind onto the NS in this
  case would be given by an expression of the form
  $\dot{M}_{\mathrm{acc}} = \dot{M}_{\mathrm{OB}} f(\tan\beta, q)$, where
  $q = M_{\mathrm{NS}}/M_{\mathrm{OB}}$ is the mass ratio of the binary and
  $\beta$ is the angle in the frame of the NS between the wind velocity
  direction and radial direction away from the OB star. Obtain the expression
  for $f(\tan\beta, q)$ assuming $M_{\mathrm{OB}} \gg M_{\mathrm{NS}}$.
  \hfill[6\,pt]
\item[(T07.4)] Consider that the fully ionized material accretes radially and
  is hindered by the strong magnetic field $\vec{B}$ of the NS. The pressure
  exerted by the magnetic field is given by $\frac{B^2}{2\mu_0}$. We shall
  assume that the NS has a dipole magnetic field whose magnitude in the
  equatorial plane varies with the distance $r$ from the NS for
  $r \gg R_{\mathrm{NS}}$ as
  \[ B(r) = B_0\left(\frac{R_{\mathrm{NS}}}{r}\right)^3 \]
  where $B_0$ is the magnetic field at the equator of the NS. Assume that the
  axis of the magnetic dipole aligns with the rotation axis of the NS.
  \begin{description}
  \item[(T07.4a)] Obtain the magnetic pressure, $P_{\mathrm{eq,\,mag}}$, in
    the equatorial plane in terms of $B_0$, $R_{\mathrm{NS}}$, $r$, and other
    suitable constants. \hfill[1\,pt]
  \item[(T07.4b)] The maximum distance where the accretion flow is stopped by
    the magnetic field at the equatorial plane is called the magnetospheric
    radius $R_{\mathrm{m}}$. This flow of matter will exert a pressure due to
    the relative motion between incoming stellar wind and the NS. Obtain an
    approximate expression for the critical magnetic field $B_{0,\mathrm{c}}$
    for which $R_{\mathrm{m}}$ coincides with $R_{\mathrm{acc}}$ and calculate
    its value in Tesla. Magnetic effects are neglected for
    $r > R_{\mathrm{m}}$ and consider $v_{\mathrm{w}} \gg v_{\mathrm{orb}}$.
    \hfill[7\,pt]
  \end{description}
\end{description}

\probhead{10.18}{Shadow of a black hole}{2025}{Mumbai, India}{TH T08}{20}
% source id: 2025_TH_T08   tag: photon sphere, BH shadow
The Event Horizon Telescope (EHT) has released an image of the supermassive
black hole at the centre of the M87 galaxy as shown in the left panel of the
figure below.

To understand some simple features of this image, we will consider a
simplified model of a non-rotating, static, spherically symmetric black hole
of mass $M = 6.5 \times 10^9\,\mathrm{M}_\odot$, surrounded by a massless,
thin, planar accretion disk of inner and outer radii,
$a_{\mathrm{inner}} = 6R_{\mathrm{SC}}$ and
$a_{\mathrm{outer}} = 10R_{\mathrm{SC}}$, respectively, where
$R_{\mathrm{SC}}$ is the Schwarzschild radius. A face-on view sketch is shown
in the right panel of the figure below (figure is not to scale).

\ioaafig{2025_TH_T08-f2.pdf}
\ioaafig{2025_TH_T08-f1.pdf}

We assume that the accretion disk is the only source of light to be
considered. Every point on the disk emits light in all directions. This light
travels under the influence of the gravitational field of the black hole. The
path of the light rays is governed by two equations given below (which are
similar to those of an object around the Sun):
\[ \frac{1}{2}v_{\mathrm{r}}^2 + \frac{L^2}{2r^2}
   \left(1 - \frac{2GM}{c^2 r}\right) = E
   \quad ; \quad
   v_\phi = r\,\omega = \frac{L}{r} \]
where $r \in (R_{\mathrm{SC}}, \infty)$ is the radial coordinate,
$\phi \in [0, 2\pi)$ is the azimuthal angle, and $E$ and $L$ are constants
related to the conserved energy and conserved angular momentum, respectively.

Here $v_{\mathrm{r}} \equiv dr/dt$ is the magnitude of the radial velocity,
$v_\phi$ is the magnitude of the tangential velocity, and
$\omega \equiv d\phi/dt$ is the angular velocity. We define the impact
parameter $b$ for a trajectory as $b = L/\sqrt{2E}$. Time dilation is
neglected in this problem.

Another useful equation is obtained by differentiating the first equation:
\[ \frac{dv_{\mathrm{r}}}{dt} - \frac{L^2}{r^3} + \frac{3GML^2}{c^2 r^4} = 0
\]

\begin{description}
\item[(T08.1)] Circular light trajectories can exist around the black hole.
  Find the radius, $r_{\mathrm{ph}}$, and impact parameter,
  $b_{\mathrm{ph}}$, for such photon trajectories in terms of $M$ and relevant
  constants. \hfill[4\,pt]
\item[(T08.2)] Calculate the time $T_{\mathrm{ph}}$ taken for completing one
  full orbit of the circular light trajectory in seconds. \hfill[2\,pt]
\item[(T08.3)] The radial velocity equation given above (the first equation in
  this question) for light trajectories can be compared with an equation of
  the form $\frac{v_{\mathrm{r}}^2}{2} + V_{\mathrm{eff}}(r) = E$. A schematic
  plot of $V_{\mathrm{eff}}/L^2$ as a function of $r$ is given below.

  \ioaafig{2025_TH_T08-f3.pdf}

  \begin{description}
  \item[(T08.3a)] The plot indicates two special radii, $r_\alpha$ and
    $r_\beta$. Obtain expressions for $r_\alpha$ and $r_\beta$ in terms of $M$
    and relevant constants. \hfill[2\,pt]
  \item[(T08.3b)] A photon travelling inward from the accretion disk towards
    the black hole can still escape out to infinity in some cases. Find the
    expression for the smallest value of the turning point radius,
    $r_{\mathrm{t}}$, for such a photon, in terms of $M$ and relevant
    constants. Find the expression for the minimum value of the impact
    parameter, $b_{\mathrm{min}}$, for this photon. \hfill[3\,pt]
  \end{description}
\item[(T08.4)] A ray of light coming from a radius $r_{\mathrm{actual}}$ from
  the centre of the system in the plane of the sky will suffer strong bending
  due to the gravity of the black hole, and eventually reach an observer
  (denoted by an eye) at a large distance $d$ from the system, as shown below.

  \ioaafig{2025_TH_T08-f4.pdf}

  To this observer, the ray would appear to have originated from a different
  point at a distance $r_{\mathrm{app}} \approx b$ from the black hole centre
  in the plane of the sky, where $b$ is the impact parameter for that photon
  trajectory. For points on the accretion disk at $r = r_{\mathrm{actual}}$,
  one may assume the following relation:
  \[ b(r_{\mathrm{actual}}) \approx r_{\mathrm{actual}}
     \left(1 + R_{\mathrm{SC}}/r_{\mathrm{actual}}\right)^{1/2} \]
  For the distant observer, like ourselves, with a face-on view of the
  accretion disk, the image of the system will appear to be circularly
  symmetric in the plane of the sky. Determine the outermost apparent radius,
  $r_{\mathrm{outer}}$, and the innermost apparent radius,
  $r_{\mathrm{inner}}$, of the image in units of au. \hfill[5\,pt]
\item[(T08.5)] Consider an isolated supermassive black hole of mass
  $M = 6.5 \times 10^9\,\mathrm{M}_\odot$ without any accretion disk. A brief
  strong burst of electromagnetic radiation occurs for 5\,s at a point Z at a
  distance, say, $r_{\mathrm{Z}} = 6R_{\mathrm{SC}}$ from the black hole as
  shown in the figure. The burst at point Z emits light in all directions. An
  observer at a point far from the black hole (denoted by an eye in the figure
  below) takes a long exposure image of the region around the black hole for
  60\,s.

  \ioaafig{2025_TH_T08-f5.pdf}

  Choose the correct option for each of the statements below:
  \begin{description}
  \item[(T08.5a)] The number of possible paths for light to travel from Z to
    the observer is\\
    (A) At most one \quad (B) Exactly one \quad (C) Exactly two \quad
    (D) Greater than two. \hfill[2\,pt]
  \item[(T08.5b)] The number of images of the EM burst at Z that will be seen
    in the long exposure image is\\
    (A) At most one \quad (B) Exactly one \quad (C) Exactly two \quad
    (D) Greater than two. \hfill[2\,pt]
  \end{description}
\end{description}

\probhead{10.19}{Hawking Radiation from Black Holes}{2025}{Mumbai, India}{TH T12}{50}
% source id: 2025_TH_T12   tag: BH temperature, evaporation, primordial BHs
\begin{description}
\item[(T12.1)] \textbf{Stellar mass blackholes}\\
  A black hole (BH) typically forms by the gravitational collapse of a massive
  star at the end of its life cycle to a point called a singularity. Due to
  the extreme gravity of such an object, nothing that enters the so-called
  event horizon (a spherical surface with $r = R_{\mathrm{SC}}$, where $r$ is
  the distance from the singularity) is able to escape from it. Here,
  $R_{\mathrm{SC}}$ is referred to as the Schwarzschild radius.
  \begin{description}
  \item[(T12.1a)] \textbf{Modelling the origin of Hawking radiation:}
    Consider a pair of particles, each with mass $m$, produced on either side
    of the BH horizon. One particle is slightly outside the horizon at
    $r \approx R_{\mathrm{SC}}$, while the other particle is inside the
    horizon at $r = \kappa R_{\mathrm{SC}}$. Assume that the total energy of a
    particle is the sum of its rest mass energy $mc^2$ and the gravitational
    potential energy due to the BH. Determine the value of $\kappa$ for which
    the particle pair has zero total energy. \hfill[4\,pt]
  \item[(T12.1b)] \textbf{Temperature of a black hole:} If the particle
    produced outside the horizon in the above process has enough kinetic
    energy, it may escape the BH in a process called Hawking radiation. The
    one inside the horizon, which has negative energy, gets absorbed and
    decreases the mass of the BH.

    Assume that all Hawking radiation is made of photons with a black body
    spectrum which peaks at the wavelength
    $\lambda_{\mathrm{bb}} \approx 16R_{\mathrm{SC}}$. It is known that for a
    solar mass BH, $R_{\mathrm{SC},\,\odot} = 2.952$\,km.

    Obtain an expression for the temperature, $T_{\mathrm{bh}}$, of the BH
    corresponding to this black body radiation, in terms of its mass
    $M_{\mathrm{bh}}$ and physical constants. Calculate the Schwarzschild
    radius, $R_{\mathrm{SC},\,10\odot}$, and temperature,
    $T_{\mathrm{bh},\,10\odot}$, for a BH with mass $10\,\mathrm{M}_\odot$.
    \hfill[4\,pt]
  \item[(T12.1c)] \textbf{Mass loss of a black hole:} Assume that the Hawking
    radiation is emitted out from the event horizon. Using the mass-energy
    equivalence, obtain an expression for the rate of mass loss,
    $dM_{\mathrm{bh}}(t)/dt$, in terms of the mass $M_{\mathrm{bh}}(t)$ of the
    BH and physical constants.

    Hence, obtain an expression for $M_{\mathrm{bh}}(t)$ for a BH with initial
    mass $M_0$. Sketch $M_{\mathrm{bh}}(t)$ as a function of $t$ from
    $M_{\mathrm{bh}} = M_0$ to $M_{\mathrm{bh}} = 0$. \hfill[8\,pt]
  \item[(T12.1d)] \textbf{Lifetime of a black hole:} Obtain an expression for
    the lifetime $\tau_{\mathrm{bh}}$ at which a black hole with initial mass
    $M_0$ completely evaporates due to Hawking radiation, in terms of $M_0$
    and physical constants. Calculate the lifetime
    $\tau_{\mathrm{bh},\,10\odot}$ (in seconds) for a black hole with
    $M_0 = 10\,\mathrm{M}_\odot$. \hfill[3\,pt]
  \item[(T12.1e)] \textbf{Black hole in a CMB radiation bath:} Consider an
    isolated black hole in space, far away from other bodies, with a current
    temperature $T_{\mathrm{bh}}^{\mathrm{now}}$, surrounded by the cosmic
    microwave background (CMB) with a current temperature
    $T_{\mathrm{cmb}}^{\mathrm{now}} = 2.7$\,K. The black hole can grow in
    mass by absorbing CMB radiation, and lose its mass by Hawking radiation.

    Taking into account the accelerating expansion of the Universe, identify
    which of the following figures show the long-term time evolution of
    $T_{\mathrm{bh}}$ in the following three cases: (X)
    $T_{\mathrm{bh}}^{\mathrm{now}} > T_{\mathrm{cmb}}^{\mathrm{now}}$, (Y)
    $T_{\mathrm{bh}}^{\mathrm{now}} = T_{\mathrm{cmb}}^{\mathrm{now}}$, (Z)
    $T_{\mathrm{bh}}^{\mathrm{now}} < T_{\mathrm{cmb}}^{\mathrm{now}}$.

    \ioaafig{2025_TH_T12-f1.pdf}

    Indicate your answer by ticking the appropriate box (only one) for each
    case X, Y or Z in the Table given in the Summary Answersheet corresponding
    to the appropriate figure number. \hfill[6\,pt]
  \end{description}
\item[(T12.2)] Primordial black holes (PBHs) of much smaller masses can form
  in the very early Universe. All the following questions are related to PBHs.
  Here, any processes that increase the mass of the black hole may be
  neglected.
  \begin{description}
  \item[(T12.2a)] \textbf{PBH evaporating at the current epoch:} As you may
    have noticed from the answers to the previous questions, black holes of
    solar mass would take a long time to evaporate. However, since PBHs can
    have a much smaller mass, we may be able to see them evaporating in
    current times.

    Find the initial mass $M_{0,\,\mathrm{PBH}}$ (in kg), Schwarzschild radius
    $R_{\mathrm{SC},\,\mathrm{PBH}}$ (in m), and temperature
    $T_{\mathrm{PBH}}$ (in K) of a black hole that may be evaporating away
    completely at the present epoch, i.e., those with lifetime
    $\tau_{\mathrm{PBH}} = 14$ billion years. \hfill[4\,pt]
  \item[(T12.2b)] \textbf{Formation of a PBH:} In the radiation-dominated
    early Universe, the scale factor varies as $a(t) \sim t^{1/2}$. In this
    era, PBHs form due to the collapse of all energy contained in a region of
    physical size $ct$, where $t$ is the age of the Universe at that time.

    A PBH with mass of $10^{12}$\,kg forms when the age of the Universe is
    about $10^{-23}$\,s. Calculate the age of the Universe, $t_{20}$, when a
    PBH of mass $10^{20}$\,kg forms. \hfill[6\,pt]
  \item[(T12.2c)] \textbf{Observed spectrum of Hawking radiation from PBH:}
    Consider a PBH of initial mass $10^{10}$\,kg which completely evaporates
    at the end of its lifetime $\tau_{\mathrm{PBH}}$. For this part, assume
    for simplicity that most of the Hawking radiation is emitted at this time,
    with a temperature corresponding to its initial mass. Also, take the scale
    factor of the Universe to be evolving as $a(t) \sim t^{2/3}$.

    Calculate the peak wavelength of this Hawking radiation as observed at
    Earth, $\lambda_{\mathrm{earth}}$, at the present epoch (at $t = 14$
    billion years). \hfill[5\,pt]
  \item[(T12.2d)] \textbf{High energy cosmic radiation from PBH:} Now assume
    that the Hawking radiation emitted at a given time $t$ corresponds to
    photons emitted with an energy $k_{\mathrm{B}}T_{\mathrm{bh}}(t)$. Also,
    the highest possible temperature for a black hole is the Planck
    temperature $T_{\mathrm{Planck}}$ where
    $k_{\mathrm{B}}T_{\mathrm{Planck}} = 10^{19}$\,GeV.

    The evolution of the scale factor over relevant time scales is given in
    the following figure. The scale factor today is set to be unity. $t(s)$ on
    the time axis represents the age of the universe in seconds.
\ioaafig{2025_TH_T12-f2.pdf}

    If a photon with an energy of $E_{\mathrm{det}} = 3.0 \times 10^{20}$\,eV
    is observed on Earth, determine the largest and the smallest possible
    values of the initial mass of the PBH ($M_0^{\mathrm{max}}$ and
    $M_0^{\mathrm{min}}$, respectively) which could be responsible for this
    photon. \hfill[10\,pt]
  \end{description}
\end{description}

\section*{Data-analysis problems}

\probhead{10.20}{Isolated black hole}{2023}{Katowice, Poland}{DA Q2}{75}
% source id: 2023_DA_Q2   tag: astrometric microlensing BH mass
In 2022, two independent groups reported the discovery of an isolated black
hole based on observations of the gravitational microlensing event
OGLE-2011-BLG-0462. In this problem, we will analyze data from the Hubble
Space Telescope to reproduce their findings.

Gravitational microlensing occurs when the light of a distant star (the
`source') is bent and magnified by the gravitational field of an intervening
object (the `lens'). The characteristic angular scale of gravitational
microlensing events, called the angular Einstein radius $\theta_E$, depends on
the mass $M$ and distance $D_\ell$ from the Earth to the lens:
\[ \theta_{\mathrm{E}} = \sqrt{\frac{4GM}{c^2}\,
   \frac{D_s - D_\ell}{D_s D_\ell}}, \]
where $D_s$ is the distance to the source star. For typical microlensing
events observed in the Milky Way, the source stars are in the Galactic bulge,
near the Galactic center, so $D_s \approx 8$\,kpc.

\begin{description}
\item[(a)] Calculate the angular Einstein radius in milliarcseconds (mas) for
  an example lens of $1\,M_\odot$ located at a distance of 1\,kpc.
  \hfill[2\,pt]
\end{description}

Suppose that at time $t$ the lens and the source are separated by an angle
$\theta \equiv u(t)\theta_E$ on the sky. Two images of the source are created
on a line through the positions of the source and the lens, at angular
distances $\theta_+$ and $\theta_-$ from the lens given by:
\[ \theta_\pm = \frac{1}{2}\left(u \pm \sqrt{u^2 + 4}\right)
   \theta_{\mathrm{E}}. \]
These two images are magnified, relative to the unlensed brightness of the
source. The absolute amplification of the images is:
\[ A_\pm = \frac{1}{2}\left(\frac{u^2 + 2}{u\sqrt{u^2 + 4}} \pm 1\right). \]
The image below shows the geometry of the event. The position of the lens is
marked as $L$, the unlensed position of the source is marked as $S$, while
$A_+$ and $A_-$ mark the positions of the two images of the source. The
dashed circle has a radius of one Einstein radius.

\ioaafig{2023_DA_Q2-f1.pdf}

\begin{description}
\item[(b)] Current telescopes cannot normally resolve this pair of images, but
  only measure the position of the image centroid, i.e.\ the
  brightness-weighted mean of the positions of the two images. Derive an
  expression for the angular separation $\theta_c$ of the image centroid
  relative to the lens as a function of $u$ and $\theta_E$. \hfill[8\,pt]
\item[(c)] Derive an expression for the source deflection $\Delta\theta$,
  i.e.\ the difference between the location of the centroid and the unlensed
  position of the source, as a function of $u$ and $\theta_{\mathrm{E}}$. What
  is the source deflection when the lens and the source are nearly perfectly
  aligned ($u \approx 0$)? \hfill[4\,pt]
\end{description}

The source and lens are moving relative to each other in the sky. Thus, both
the total magnification of the images and the position of the centroid changes
with time, resulting in observable photometric and astrometric microlensing
effects. For now, we assume that the source-lens relative motion is
rectilinear.

The plot below shows the light curve of the gravitational microlensing event
OGLE-2011-BLG-0462, discovered by the OGLE sky survey led by astronomers from
the University of Warsaw. The solid line shows the best-fitting light curve
model. The Einstein timescale of the event, i.e.\ the time needed for the
source to move by one angular Einstein radius relative to the lens, was
$t_E = 247$ days. The event peaked on 21 July 2011 (HJD = 2455763). The
minimal separation between the lens and the source was $u_0 \approx 0$.

\ioaafig{2023_DA_Q2-f2.pdf}

The table below shows the measured positions of the source star against the
background objects in the East and North directions based on images from the
Hubble Space Telescope.

\begin{center}
\begin{tabular}{rrr}
\toprule
HJD & E position (mas) & N position (mas) \\
\midrule
2455765.2 & $2.58 \pm 0.13$ & $7.29 \pm 0.16$ \\
2455865.7 & $2.32 \pm 0.12$ & $5.44 \pm 0.24$ \\
2456179.7 & $0.46 \pm 0.14$ & $1.62 \pm 0.08$ \\
2456195.8 & $0.88 \pm 0.36$ & $1.56 \pm 0.77$ \\
2456426.2 & $-1.02 \pm 0.21$ & $-0.94 \pm 0.12$ \\
2456587.7 & $-2.04 \pm 0.07$ & $-1.88 \pm 0.40$ \\
2456956.6 & $-4.54 \pm 0.25$ & $-5.16 \pm 0.29$ \\
2457995.2 & $-11.14 \pm 0.12$ & $-15.14 \pm 0.17$ \\
\bottomrule
\end{tabular}
\end{center}

\begin{description}
\item[(d)] Plot the measured positions of the source star in the East and
  North directions as a function of time. \hfill[10\,pt]
\item[(e)] The observed motion of the source star is the sum of two effects:
  rectilinear proper motion of the source and astrometric microlensing
  effects. Calculate the proper motion (in mas/year) of the source in the East
  and North directions and its uncertainty. \hfill[8\,pt]
\item[(f)] After subtracting the effects of proper motion from the data,
  calculate and plot the total resultant astrometric deflection as a function
  of $u$. Neglect uncertainty of the proper motion determination.
  \hfill[20\,pt]
\item[(g)] Analyse the data to determine the angular Einstein radius
  $\theta_{\mathrm{E}}$ of the event and its uncertainty. (\textit{Hint:} it
  may be helpful to linearise the expression for $\Delta\theta$).
  \hfill[16\,pt]
\item[(h)] For long-timescale events such as OGLE-2011-BLG-0462, the
  rectilinear approximation of the relative lens-source proper motion is not
  strictly true and the orbital motion of the Earth has to be taken into
  account. This allows measurement of a dimensionless quantity called the
  microlensing parallax, defined as
  $\pi_E = (\pi_l - \pi_s)/\theta_E$, where $\pi_l$ and $\pi_s$ are parallaxes
  of the lens and the source, respectively.

  For this event $\pi_E = 0.095 \pm 0.009$. Rearrange the expression for
  $\theta_{\mathrm{E}}$ given earlier to calculate the mass of the lens in
  solar masses and its uncertainty. \hfill[7\,pt]
\end{description}

\ioaafig{2023_DA_Q2-f3.pdf}
\ioaafig{2023_DA_Q2-f4.pdf}

\clearpage
\section*{Solutions to Chapter 10}

\solhead{10.1}{}
% source id: 2007_TH_1.14
The spherical boundary at which the escape velocity becomes equal to the speed
of light is of radius $R = \dfrac{2GM_\odot}{c^2}$
\begin{align*}
R &= \frac{2 \times 6.672 \times 10^{-11} \times 1.989 \times 10^{30}}
         {\left(2.99792458 \times 10^{8}\right)^2}~\mathrm{m}\\
  &= \frac{2 \times 6.672 \times 1.989}{2.998 \times 2.998} \times 10^{3}~\mathrm{m}\\
  &= 2.95~\mathrm{km}
\end{align*}

\solhead{10.2}{Gravitational Lensing}
% source id: 2007_TH_Q4
\ioaafig{2007_TH_Q4-s1.pdf}
\begin{description}
\item[a).] From figure 1, for small angles, $\tan\theta \approx \theta$, hence
\begin{align*}
\mathrm{PS}_1 &= \mathrm{PS} + \mathrm{SS}_1\\
\theta D_S &= \beta D_S + \left(D_S - D_L\right)\phi\\
\theta - \beta &= \frac{4GM}{\xi c^2}\,\frac{\left(D_S - D_L\right)}{D_S}
  \tag*{(1)\quad(1 point)}
\end{align*}
Note that $\theta = \dfrac{\xi}{D_L}$ also, hence,
\[
\theta^2 - \beta\theta
  = \left(\frac{4GM}{c^2}\right)\left(\frac{D_S - D_L}{D_L D_S}\right)
  \tag*{(2)}
\]
For a perfect alignment in which $\beta = 0$, we have $\theta = \pm\theta_E$,
where
\[
\theta_E = \sqrt{\left(\frac{4GM}{c^2}\right)\left(\frac{D_S - D_L}{D_L D_S}\right)}
  \tag*{(3)\quad(1 point)}
\]

\item[b).] From equation (3), for a solar-mass lens with $D_S = 50$~kpc,
$D_L = 50 - 10 = 40$~kpc,
\begin{align*}
\theta_E &= \sqrt{\left(\frac{4GM}{c^2}\right)
                  \left(\frac{D_S - D_L}{D_L D_S}\right)}
          = 0.956 \times 10^{-9}~\text{radian}\\
         &= 1.97 \times 10^{-4}~\text{arc second} \tag*{(1 point)}
\end{align*}

\item[c.)] The resolution of the Hubble space telescope having diameter of
2.4~m is
\[
\theta_{Hubble} = 1.22\,\frac{\lambda}{D}
  = 1.22 \times \frac{5 \times 10^{-7}~\mathrm{m}}{2.4~\mathrm{m}}
  = 2.54 \times 10^{-7}~\text{radian}
\]
for light of wavelength 500~nm. \hfill(1 point)

Hence the Hubble telescope could not resolve \underline{this} Einstein ring.
\hfill(1 point)

\item[d).] The quadratic equation (2) has two distinct roots, namely,
\begin{align*}
\theta_1 &= \frac{\beta}{2}
  + \sqrt{\left(\frac{\beta}{2}\right)^2 + \theta_E^2} \tag*{(1 point)}\\
\theta_2 &= \frac{\beta}{2}
  - \sqrt{\left(\frac{\beta}{2}\right)^2 + \theta_E^2} \tag*{(1 point)}
\end{align*}
where $\theta_E = \sqrt{\left(\dfrac{4GM}{c^2}\right)
\left(\dfrac{D_S - D_L}{D_L D_S}\right)}$.

This implies that there are two images for a single isolated source.

\item[e).]
\begin{align*}
\frac{\theta_{1,2}}{\beta}
  &= \frac{\dfrac{\beta}{2}
       \pm \sqrt{\left(\dfrac{\beta}{2}\right)^2 + \theta_E^2}}{\beta}\\
  &= \frac{1}{2}
     \pm \sqrt{\left(\frac{1}{2}\right)^2 + \left(\frac{1}{\eta}\right)^2}
   = \frac{1}{2}\left[1 \pm \frac{\sqrt{\eta^2 + 4}}{\eta}\right]
  \tag*{(1 point)}
\end{align*}

\item[f).] From equation (2) $\theta^2 - \beta\theta - \theta_E^2 = 0$
\begin{align*}
\left(\theta + \Delta\theta\right)^2
  - \left(\beta + \Delta\beta\right)\left(\theta + \Delta\theta\right)
  - \theta_E^2 &= 0\\
\frac{\Delta\theta}{\Delta\beta}
  &= \frac{\theta}{2\theta - \beta} \tag*{(1 point)}\\
\left(\frac{\Delta\theta}{\Delta\beta}\right)_{\theta=\theta_{1,2}}
  = \frac{\theta_{1,2}}{2\theta_{1,2} - \beta}
  &= \frac{\dfrac{1}{2}\eta \pm \sqrt{1 + \dfrac{\eta^2}{4}}}
          {\pm 2\sqrt{1 + \dfrac{\eta^2}{4}}}
   = \frac{1}{2}\left[1 \pm \frac{\eta}{\sqrt{\eta^2 + 4}}\right]
  \tag*{(1 point)}
\end{align*}
\end{description}

\solhead{10.3}{}
% source id: 2009_TH_P1
Schwarzschild radius of a black hole with mass $M$ is
\[
R = \frac{2GM}{c^2} \tag*{(4 points)}
\]
Then the mass density can be estimated as
\begin{align*}
\rho &= \frac{M}{\frac{4}{3}\pi \frac{8G^3M^3}{c^6}}
      = \frac{3c^6}{32\pi}\,\frac{1}{G^3M^2} \tag*{(2 points)}\\
     &= \frac{3 \times (3.00 \times 10^{8})^6}{32 \times 3.1416}\,
        \frac{1}{(6.67 \times 10^{-11})^3 (10^{8} \times 1.99 \times 10^{30})^2}\\
     &= 1.85 \times 10^{3}~\mathrm{kg\,m^{-3}} \tag*{(4 points)}
\end{align*}

\solhead{10.4}{}
% source id: 2009_TH_P10
Because of high conductivity of plasma inside the star, flux of magnetic field
will be conserved through contraction, then:
\[
4\pi R^2 B = 4\pi R_n^2 B_n \tag*{(6 points)}
\]
Where $R_n$ and $B_n$ are radius and magnetic field of neutron star, thus:
\begin{align*}
B_n &= \left(\frac{R}{R_n}\right)^2 B
     = \left(\frac{4 \times 6.96 \times 10^{5}}{20}\right)^2 \tag*{(2 points)}\\
B_n &= 1.93 \times 10^{10}~\mathrm{Gauss}\\
    &= 1.93 \times 10^{6}~\mathrm{T} \tag*{(2 points)}
\end{align*}

\solhead{10.5}{}
% source id: 2010_TH_S2
\begin{align*}
\sqrt{\frac{2GM_{object}}{R_{object}}} &> c \tag*{(4 points)}\\
R_{object} &< \frac{2GM_{object}}{c^2} \tag*{(2 points)}\\
R_{object} &< \frac{2 \times 6.6726 \times 10^{-11} \times 1.9891 \times 10^{30}}
                  {(2.9979 \times 10^{8})^2}\\
R &< 2953.6~\mathrm{m} \tag*{(4 points)}
\end{align*}

\solhead{10.6}{}
% source id: 2012_TH_S8
As for the probability, supposing that the neutron star space orientation is
random, is
\[
p = \frac{2A_p(d)}{4\pi d^2} \tag*{(2 pt)}
\]
where $A_p(d)$ is the area of the path of one light cone.

\ioaafig{2012_TH_S8-s1.pdf}

As the figure above shows, the area is the difference of the spherical caps of
height $H$ and $D$:
\[
A = 2\pi d (H - D)
  = 2\pi d^2 \left[\cos\left(\theta - \frac{\alpha}{2}\right)
    - \cos\left(\theta + \frac{\alpha}{2}\right)\right] \tag*{(3 pt)}
\]
\[
p = \cos\left(\theta - \frac{\alpha}{2}\right)
  - \cos\left(\theta + \frac{\alpha}{2}\right)
  = 3.5 \cdot 10^{-2} \quad \text{or} \quad 3.5\% \tag*{(2 pt)}
\]
The light waves are spread along the double cone and not along a sphere, thus:
\[
F = \frac{L}{2A(d)}
  = \frac{10\,000\,L_\odot}{4\pi d^2\left(1 - \cos\dfrac{\alpha}{2}\right)}
\]
Where $A(d)$ is the area of one light cone. Using the Sun's absolute magnitude,
we can find the absolute magnitude of the pulsar ($m_p$):
\[
m_p - M_\odot
  = -2.5\log\left(\frac{F}{\dfrac{L_\odot}{4\pi(10~\mathrm{pc})^2}}\right)
  = -2.5\log\left(10\,000 \cdot \frac{1}{1 - \cos\dfrac{\alpha}{2}}
    \cdot \left(\frac{10}{1000}\right)^2\right) \tag*{(3 pt)}
\]
\[
m_p = -3.24
\]

\solhead{10.7}{Cosmic radiation}
% source id: 2014_TH_P3
\begin{description}
\item[(A)] In the disintegration process the laws of energy conservation and
the law of the conservation of momentum are both obeyed.

In the general case the law of conservation of the momentum is represented in
the down below figure.
\ioaafig{2014_TH_P3-s1.pdf}

The total initial energy of the $\pi^0$ meson is
\[
E^2 = p^2c^2 + m_0^2c^4
\]
And its kinetic energy is
\[
E_c = E - m_0 c^2
\]
The expressions of the 2 conservation laws written after the disintegration
are:
\begin{align*}
\vec{p} &= \vec{p}_1 + \vec{p}_2;\\
E + m_0 c^2 &= E_1 + E_2,
\end{align*}
The energy of the photon 1 can be calculated using the notations in the figure
\begin{align*}
E &= E_c + m_0 c^2;\\
p_1\sin\theta_1 &= p_2\sin\theta_2;\\
\frac{E_1}{c}\sin\theta_1 &= \frac{E_2}{c}\sin\theta_2;\\
E^2 &= p^2c^2 + m_0^2c^4;\\
p^2c^2 &= E^2 - m_0^2c^4;\\
E &= E_1 + E_2;\\
E_2 &= E - E_1 = \left(E_c + m_0c^2\right) - E_1;\\
E_1 &= \frac{m_0^{\,2}c^4}{2}\,
  \frac{1}{E_c + m_0c^2 - \cos\theta_1\sqrt{E_c\left(E_c + m_0c^2\right)}}.
\end{align*}
Similar the second photon energy is:
\[
E_2 = \frac{m_0^{\,2}c^4}{2}\,
  \frac{1}{E_c + m_0c^2 - \cos\theta_2\sqrt{E_c\left(E_c + m_0c^2\right)}}
\]

\item[(B)] If one of the photons has the maximum possible energy
$E_{\mathrm{max}}$ and consequently the other photon has the minimum possible
energy $E_{\mathrm{min}}$ the law of momentum conservation is sketched:
\ioaafig{2014_TH_P3-s2.pdf}

Thus the relations become very simple:
\begin{align*}
mv &= \frac{m_0 v}{\sqrt{1 - \dfrac{v^2}{c^2}}}\\
E_{\mathrm{max}} &= \frac{m_0 c^2}{2\sqrt{1 - \dfrac{v^2}{c^2}}}
  \left(1 + \frac{v}{c}\right);\\
E_{\mathrm{min}} &= \frac{m_0 c^2}{2\sqrt{1 - \dfrac{v^2}{c^2}}}
  \left(1 - \frac{v}{c}\right);
\end{align*}

\item[(C)]
\[
v = c\,\frac{E_{\mathrm{max}} - E_{\mathrm{min}}}
            {E_{\mathrm{max}} + E_{\mathrm{min}}}.
\]
\end{description}

\solhead{10.8}{}
% source id: 2015_TH_S12
\begin{align*}
E_0 &= m_0 c^2 = 941~\mathrm{MeV} \tag*{(20 points)}\\
\gamma &= \frac{1}{\sqrt{1 - \dfrac{v^2}{c^2}}}
  = \frac{K}{E_0} + 1
  = \frac{1000~\mathrm{MeV}}{941~\mathrm{MeV}} + 1 = 2.062 \tag*{(20 points)}\\
1 - \frac{v^2}{c^2} &\approx \frac{1}{4}
  \;\Rightarrow\; v \approx 0.87c \tag*{(30 points)}\\
\text{Elapse time} &= \frac{\text{Earth--Sun distance}}{\text{speed}}
  = \frac{1.50 \times 10^{11}~\mathrm{m}}
         {0.87 \times (3 \times 10^{8})~\mathrm{m/s}}\\
  &= 9.6~\text{minutes} \tag*{(30 points)}
\end{align*}

\solhead{10.9}{Gravitational Lensing Telescope}
% source id: 2016_TH_T10
\begin{description}
\item[(T10.1)] The rays travelling closer to the gravitational body will bend
more. Thus, we get shortest convergence point where the rays just grazing the
solar surface will meet each other.
\ioaafig{2016_TH_T10-s1.pdf}
\begin{align*}
\theta_{\mathrm{b}} &= \frac{2R_{\mathrm{sch}}}{R_\odot}
  \simeq \frac{R_\odot}{f_{\min}} \tag*{(2.0)}\\
\therefore f_{\min} &= \frac{R_\odot^2}{2R_{\mathrm{sch}}}
  = \frac{R_\odot^2 c^2}{4GM_\odot} \tag*{(2.0)}\\
  &= \frac{(6.955 \times 10^{8} \times 2.998 \times 10^{8})^2}
          {4 \times 6.674 \times 10^{-11} \times 1.989 \times 10^{30}}~\mathrm{m}\\
  &= 8.188 \times 10^{13}~\mathrm{m}
   = \frac{8.188 \times 10^{13}}{1.496 \times 10^{11}}~\mathrm{AU}
\end{align*}
\[
\boxed{f_{\min} = 547.3~\mathrm{AU}} \tag*{(2.0)}
\]

\item[(T10.2)] The following figure needs to be drawn. \hfill(1.0)
\ioaafig{2016_TH_T10-s2.pdf}

The light bending from the surface of the Sun ($r = R_\odot$) will intersect
the detector at its centre, as it is kept at $f_{\min}$.

The boundary of the detector will be intersected by a light ray with
$r = R_\odot + h$. This ray will intersect the $x$-axis at a distance $f_2$.
\[
f_2 = \frac{(R_\odot + h)^2}{2R_{\mathrm{sch}}} \tag*{(2.0)}
\]
\textbf{Same argument as in Part 1.} For small angles,
\begin{align*}
a &= (f_2 - f_{\min})\theta_2\\
  &= \left[\frac{(R_\odot + h)^2}{2R_{\mathrm{sch}}}
     - \frac{R_\odot^2}{2R_{\mathrm{sch}}}\right]
     \frac{2R_{\mathrm{sch}}}{(R_\odot + h)}
   = \frac{2R_\odot h + h^2}{R_\odot + h}\\
  &\simeq 2h \tag*{(2.0)}
\end{align*}
Let original intensity of the incoming radiation be $I_0$.

The flux at the detector in the presence of Sun is
$\Phi_\odot = I_0\, 2\pi\, R_\odot\, h$ \hfill(1.0)

The flux at the detector in the absence of Sun is
$\Phi_0 = I_0\, \pi\, a^2$ \hfill(1.0)

The amplification is therefore
\[
A_{\mathrm{m}} = \frac{\Phi_\odot}{\Phi_0}
  = \frac{I_0 2\pi R_\odot h}{I_0 \pi a^2}
  = \boxed{\frac{R_\odot}{a}} \tag*{(1.0)}
\]

\item[(T10.3)]
\ioaafig{2016_TH_T10-s3.pdf}
All rays should focus at the same spot. This should be evident from figure
drawn on answersheet or otherwise. \hfill(2.0)

Let there be two rays with impact parameters $r_1$ and $r_2$. The
corresponding distances of focus will be
\[
f_i = \frac{r_i^2}{2r_{\mathrm{sch}_i}} = \frac{r_i^2 c^2}{4GM(r_i)} \tag*{(2.0)}
\]
\textbf{Same argument as in Part 1.} The rquirement % sic: "rquirement" in source
$f_1 = f_2$ implies
\[
\frac{r_1^2}{r_2^2} = \frac{M(r_1)}{M(r_2)} \tag*{(1.0)}
\]
The required mass distribution is: $\boxed{M(r) \propto r^2}$ \hfill(1.0)
\end{description}

\solhead{10.10}{Deflection of radio photons in the gravitational field of solar system bodies}
% source id: 2019_TH_2
One has to notice that in the case of the Sun $d = R_\odot$, that is $d$
equals the solar radius. \hfill(1 p)

The angular deflection scales as $m/r$ (mass over radius) of the deflector
body, so one has to consider $M_{\mathrm{b}}/R_{\mathrm{b}}$, where b stands
for Jupiter and Moon, respectively.
\[
\frac{\Delta\theta_{\mathrm{b}}}{\Delta\theta_\odot}
  = \frac{M_{\mathrm{b}}/R_{\mathrm{b}}}{M_\odot/R_\odot}
  \tag*{(2.1)\quad(2 p)}
\]
This means that:
\[
\Delta\theta_{\mathrm{b}}
  = \Delta\theta_\odot\,\frac{M_{\mathrm{b}}/R_{\mathrm{b}}}{M_\odot/R_\odot}
  \tag*{(2.2)\quad(1 p)}
\]
\begin{description}
\item[a)] Jupiter

Substituting the masses and radii one gets:
\[
\Delta\theta_{\mathrm{J}} = \boxed{0.017''} = \boxed{17~\mathrm{mas}}
  \tag*{(2 p)}
\]
The deflection of radio waves by Jupiter's gravitational field is
\emph{greater} than 0.1\,mas, so it is \emph{observable} with the VLBI
network. The answer is: YES \hfill(1 p)

\item[b)] Moon

Substituting the masses and radii one gets:
\[
\Delta\theta_{\mathrm{M}} = \boxed{0.000\,026''} = \boxed{0.026~\mathrm{mas}}
  \tag*{(2 p)}
\]
The deflection of radio waves by Moon's gravitational field is \emph{smaller}
than 0.1\,mas, so it is \emph{not observable} with the VLBI network. The
answer is: NO \hfill(1 p)
\end{description}

\solhead{10.11}{The supermassive black hole in the centre of Milky Way Galaxy and M87}
% source id: 2019_TH_3
\begin{description}
\item[a)] The scale of the black hole shadow is set by the Schwarzschild
  radius $R_{\mathrm{S}} = 2GM/c^2$. The light ring radius is
  $3R_{\mathrm{S}}$. \hfill(2\,p)

  For a diffraction limited instrument the angular resolution is bounded by
  \begin{equation*}
    \theta_{\mathrm{res}} = \frac{2R}{d} \ge \frac{1.22\lambda}{D}
    \tag{3.1}
  \end{equation*}
  \hfill(2\,p)

  where $\lambda$ is the wavelength and $D$ is the diameter of the
  instrument.

  Using the expression for the Schwarzschild radius:
  \begin{equation*}
    D \ge \frac{1.22\lambda}{6R_{\mathrm{S}}}\,d
      = \frac{1.22c^2\lambda d}{12GM}
    \tag{3.2}
  \end{equation*}
  \hfill(2\,p)
\item[b)] For $\lambda = 1.3$\,mm, this gives
  \begin{description}
  \item[(a)] $D > 7200\,\mathrm{km} \approx 1.1\,R_\oplus$ \hfill(1\,p)
  \item[(b)] $D > 6400\,\mathrm{km} \approx 1\,R_\oplus$ \hfill(1\,p)
  \end{description}
\item[c)] (B) Interferometry with array of radio telescopes \hfill(2\,p)
\end{description}

\solhead{10.12}{Planck's units}
% source id: 2022_TH_1
% The official solution labels its single part "(a)".
(a) From dimensional analysis
\begin{align*}
  c &\equiv \left[L^1 M^0 T^{-1}\right]\\
  G &\equiv \left[L^3 M^{-1} T^{-2}\right]\\
  \hbar &\equiv \left[L^2 M^1 T^{-1}\right]\\
  k_B &\equiv \left[L^2 M^1 T^{-2} \Theta^{-1}\right]
\end{align*}
\hfill(1 pt each)

\begin{align*}
  G\hbar &\equiv \left[L^5 M^0 T^{-3}\right]\\
  \frac{G\hbar}{c^3} &\equiv \left[L^2 M^0 T^0\right]\\
  \therefore\ L_p &= \sqrt{\frac{G\hbar}{c^3}}
\end{align*}
\hfill(1.5 pt)

\begin{align*}
  \frac{\hbar c}{G} &\equiv \left[L^0 M^2 T^0\right]\\
  \therefore\ m_p &= \sqrt{\frac{\hbar c}{G}}
\end{align*}
\hfill(1.5 pt)

\begin{align*}
  \frac{L_p}{c} &\equiv \left[L^0 M^0 T^1\right]\\
  \therefore\ t_p &= \sqrt{\frac{\hbar G}{c^5}}
\end{align*}
\hfill(1.0 pt)

\begin{align*}
  T_p &= \frac{L_p^2 m_p}{k_B t_p^2}
       = \frac{G\hbar}{k_B c^3} \times \frac{c^5}{\hbar G}
         \times \sqrt{\frac{\hbar c}{G}}\\
  \therefore\ T_p &= \sqrt{\frac{\hbar c^5}{G k_B^2}}
\end{align*}
\hfill(1.0 pt)

\begin{align*}
  4\pi\epsilon_0 &\equiv \left[L^{-3} M^{-1} T^2 Q^2\right]\\
  \hbar c &\equiv \left[L^3 M^1 T^{-2}\right]\\
  \therefore\ q_p &= \sqrt{4\pi\epsilon_0 \hbar c}
\end{align*}
\hfill(1.0 pt)

\solhead{10.13}{Accretion}
% source id: 2022_TH_10
\begin{description}
\item[(a)] Consider a particle at a distance $r$ from the center of the
  compact object. In case of maximal luminosity, the gravitational force must
  be balanced by the radiation pressure.
  \[ F_G = G\frac{m_H M}{r^2} \]
  \hfill(0.5 pt)
  \[ F_R = \frac{\sigma_e L_E}{4\pi c r^2} \]
  \hfill(0.5 pt)
  \[ \text{But, } F_G = F_R \]
  \hfill(1.0 pt)
  \[ \therefore\ L_E = \frac{4\pi G m_H M c}{\sigma_e} \]
  \hfill(1.0 pt)
\item[(b)] The lower bound of its mass given by:
  \begin{align*}
    \sigma_e &= \frac{8\pi}{3}
      \left(\frac{e^2}{4\pi\varepsilon_0 m_e c^2}\right)^2\\
    &\approx 6.65 \times 10^{-29}\ \mathrm{m}^2
  \end{align*}
  \hfill(1.5 pt)
  \begin{align*}
    M &= \frac{\sigma_e L_\odot}{4\pi G m_H c}\\
      &= 6.04 \times 10^{25}\ \mathrm{kg}\\
      &= 3.04 \times 10^{-5}\,M_\odot
  \end{align*}
  \hfill(2.5 pt)
\item[(c)] Gravitational energy of atoms of total mass $\Delta m$:
  \begin{align*}
    \Delta E &= \frac{GM\Delta m}{R}\\
    L_{\mathrm{acc}} &= \frac{\Delta E}{\Delta t}\\
    \therefore\ L_{\mathrm{acc}} &= \frac{GM}{R}\frac{\Delta m}{\Delta t}
      = \frac{GM}{R}\dot{M}
  \end{align*}
  \hfill(2.0 pt)
\item[(d)] For maximum $\dot{M}$:
  \begin{align*}
    L_{\mathrm{acc}} &= L_E\\
    \frac{GM}{R}\dot{M} &= \frac{4\pi G m_H M c}{\sigma_e}\\
    \dot{M} &= \left(\frac{4\pi m_H c}{\sigma_e}\right) R
  \end{align*}
  \hfill(2.0 pt)
\item[(e)] For given value of $R_*$:
  \begin{align*}
    \dot{M} &= \left(\frac{4\pi m_H c}{\sigma_e}\right) R_*\\
    \dot{M} &= 1.14 \times 10^{21}\ \mathrm{kg/s}\\
      &= 3.59 \times 10^{28}\ \mathrm{kg/yr}\\
    \dot{M} &= 0.018\,\mathrm{M}_\odot/\mathrm{yr}
  \end{align*}
  \hfill(2.0 pt)
\item[(f)] The total angular momentum and total mass must be conserved.
  \[ L = M_1 r_1^2 \Omega + M_2 r_2^2 \Omega \]
  \hfill(1.0 pt)
  \begin{align*}
    \text{Let } M_T &= M_1 + M_2\\
    \text{and } a &= r_1 + r_2\\
    r_1 = \frac{M_2 a}{M_T} &\quad\text{and}\quad r_2 = \frac{M_1 a}{M_T}
  \end{align*}
  \hfill(1.0 pt)
  \begin{align*}
    \therefore\ L &= \left[M_1\left(\frac{M_2 a}{M_T}\right)^2
      + M_2\left(\frac{M_1 a}{M_T}\right)^2\right]\Omega\\
    &= \left(M_1 M_2^2 + M_2 M_1^2\right)\frac{a^2\Omega}{M_T^2}\\
    \therefore\ L &= \frac{a^2\Omega M_1 M_2}{M_T}
      = \frac{a^2\Omega M_1 (M_T - M_1)}{M_T}
  \end{align*}
  \hfill(1.0 pt)
  \begin{align*}
    \therefore\ a &= \left(\sqrt{\frac{L M_T}{\Omega}}\right)
      M_1^{-0.5}(M_T - M_1)^{-0.5}\\
    a &= K M_1^{-0.5}(M_T - M_1)^{-0.5}
  \end{align*}
  \hfill(1.0 pt)

  Now let us say mass of the compact object \textbf{increases} by small
  fraction $\Delta M_1$. Let the corresponding \textbf{increases} in the
  separation be $\Delta a$.
  \begin{align*}
    a + \Delta a &= K(M_1 + \Delta M_1)^{-0.5}
      (M_T - M_1 - \Delta M_1)^{-0.5}\\
    &= K M_1^{-0.5}(M_T - M_1)^{-0.5}
      \left(1 + \frac{\Delta M_1}{M_1}\right)^{-0.5}
      \left(1 - \frac{\Delta M_1}{M_T - M_1}\right)^{-0.5}\\
    a + \Delta a &= a\left(1 + \frac{\Delta M_1}{M_1}\right)^{-0.5}
      \left(1 - \frac{\Delta M_1}{M_2}\right)^{-0.5}
  \end{align*}
  \hfill(1.0 pt)
  \begin{align*}
    \therefore\ 1 + \frac{\Delta a}{a}
      &\approx \left(1 - \frac{\Delta M_1}{2M_1}\right)
        \left(1 + \frac{\Delta M_1}{2M_2}\right)\\
    &\approx 1 - \frac{\Delta M_1}{2M_1} + \frac{\Delta M_1}{2M_2}\\
    \therefore\ \frac{\Delta a}{a}
      &= \frac{\Delta M_1}{2}\left(\frac{M_1 - M_2}{M_1 M_2}\right)
  \end{align*}
  \hfill(1.0 pt)

  Thus, $\Delta a$ is positive (separation increases) when $M_1 > M_2$ and
  the separation decreased when $M_1 < M_2$. \hfill(1.0 pt)
\end{description}

\solhead{10.14}{Relativistic Beaming}
% source id: 2022_TH_13
\begin{description}
\item[(a)] Energy of a photon:
  \[ E = hf \]
  while the momentum:
  \[ \vec{p} = \frac{hf}{c}\vec{n} \]
  \hfill(2.0 pt)

  where $\vec{n}$ is an unit vector in the direction of the motion of the
  photon. Substituting this into energy-momentum transformation law we have:
  \begin{align*}
    \frac{hf_L}{c} &= \gamma\left(\frac{hf_R}{c} + p_{x_R}\frac{v}{c}\right)\\
    p_{x_L} &= \gamma\left(p_{x_R} + \frac{hf_R v}{c^2}\right)\\
    p_{y_L} &= p_{y_R}\\
    p_{z_L} &= p_{z_R}
  \end{align*}
  \hfill(2.0 pt)

  in our case:
  \begin{align*}
    p_{x_L} &= |p_L|\cos\theta_L, & p_{y_L} &= |p_L|\sin\theta_L,
      & p_{z_L} &= 0\\
    p_{x_R} &= |p_R|\cos\theta_R, & p_{y_R} &= |p_R|\sin\theta_R,
      & p_{z_R} &= 0
  \end{align*}
  \hfill(3.0 pt)

  where
  \[ |p| = \frac{hf}{c} \]
  Thus,
  \begin{align*}
    f_L &= \gamma\left(f_R + \frac{v}{h}|p_R|\cos\theta_R\right)\\
        &= \gamma\left(f_R + \frac{v}{h}\,\frac{hf_R}{c}\cos\theta_R\right)\\
    \therefore\ f_L &= \gamma f_R\left(1 + \frac{v}{c}\cos\theta_R\right)
  \end{align*}
  \hfill(2.0 pt)
  \begin{align*}
    \cos(\theta_L) &= \frac{p_{x_L}}{|p_L|} = \frac{c\,p_{x_L}}{hf_L}
      = \frac{c\gamma\left(p_{x_R} + \dfrac{hf_R v}{c^2}\right)}
             {h\gamma f_R\left(1 + \dfrac{v}{c}\cos\theta_R\right)}\\
    &= \frac{c\left(\dfrac{hf_R}{c}\cos\theta_R + \dfrac{hf_R v}{c^2}\right)}
            {hf_R\left(1 + \dfrac{v}{c}\cos\theta_R\right)}\\
    \therefore\ \cos\theta_L
      &= \frac{\cos\theta_R + \dfrac{v}{c}}
              {1 + \dfrac{v}{c}\cos\theta_R}
  \end{align*}
  \hfill(2.0 pt)
\item[(b)] Now just by plugging in the values of cosine we see that in case
  (i) (orange arrow) $\cos\theta_L = 1$, So the photon keeps moving in the
  same direction, in case (ii) (Green arrow) $\cos\theta_L = 0$, in case
  (iii) (Yellow arrow) $\cos\theta_L = v/c$, (iv) (Grey arrow)
  $\cos\theta_L = -1$.
  \ioaafig{2022_TH_13-s1.pdf}
\item[(c)] From the figure we see that:
  \[ KP = OP - OK\sin\theta = R - r\sin\theta \]
  \hfill(1.0 pt)

  where $\theta = \omega t$.
  \ioaafig{2022_TH_13-s2.pdf}

  Let $t_L = T$ be the time at which photon reaches $P$ after leaving $K$ at
  a time $t$ (in lab frame) then:
  \begin{align*}
    c(T - t) &= KP = R - r\sin\theta \approx R\\
    \therefore\ t &= t_L - \frac{R}{c}
  \end{align*}
  \hfill(1.0 pt)

  This photon makes angle $\theta = \omega t$ to the direction of its motion
  in lab frame. Now,
  \begin{align*}
    \cos\theta_L &= \frac{\cos\theta_R + \dfrac{v}{c}}
                        {1 + \dfrac{v}{c}\cos\theta_R}\\
    \cos\theta_L + \frac{v}{c}\cos\theta_R\cos\theta_L
      &= \cos\theta_R + \frac{v}{c}\\
    \cos\theta_R\left(1 - \frac{v}{c}\cos\theta_L\right)
      &= \cos\theta_L - \frac{v}{c}\\
    \therefore\ \cos\theta_R
      &= \frac{\cos\theta_L - \dfrac{v}{c}}
              {1 - \dfrac{v}{c}\cos\theta_L}
  \end{align*}
  \hfill(2.0 pt)
  \begin{align*}
    f_L &= \gamma f_R\left(1 + \frac{v}{c}\cos\theta_R\right)\\
    &= \gamma f_R\left(1 + \frac{\dfrac{v}{c}
        \left(\cos\theta_L - \dfrac{v}{c}\right)}
        {1 - \dfrac{v}{c}\cos\theta_L}\right)\\
    &= \gamma f_R\left(\frac{1 - \dfrac{v}{c}\cos\theta_L
        + \dfrac{v}{c}\cos\theta_L - \left(\dfrac{v}{c}\right)^2}
        {1 - \dfrac{v}{c}\cos\theta_L}\right)\\
    &= \frac{f_R}{\gamma\left(1 - \dfrac{v}{c}\cos\theta_L\right)}
     = \frac{f_R}{\gamma\left(1 - \dfrac{v}{c}\cos(\omega t)\right)}\\
    \therefore\ f_L
      &= \frac{f_R}{\gamma\left(1 - \dfrac{v}{c}
         \cos\left[\omega\left(t_L - \dfrac{R}{c}\right)\right]\right)}
  \end{align*}
  \hfill(3.0 pt)
\item[(d)] \hspace{0pt}
  \begin{align*}
    \Delta\Omega_R &= -\Delta(\cos\theta_R)\cdot\Delta\phi\\
    &= \left[\cos\theta_R - \cos(\theta_R + \Delta\theta_R)\right]
       \cdot\Delta\phi\\
    &= \left[\cos\theta_R - \cos\theta_R\cos(\Delta\theta_R)
       + \sin\theta_R\sin(\Delta\theta_R)\right]\cdot\Delta\phi\\
    &= \left[\cos\theta_R - \cos\theta_R\cdot(1)
       + \sin\theta_R(\Delta\theta_R)\right]\cdot\Delta\phi\\
    \Delta\Omega_R &= \sin\theta_R(\Delta\theta_R)(\Delta\phi)\\
    \therefore\ \Delta\Omega_L &= \sin\theta_L(\Delta\theta_L)(\Delta\phi)
  \end{align*}
  \hfill(3.0 pt)
  \begin{align*}
    \sin^2\theta_L &= 1 - \cos^2\theta_L\\
    &= 1 - \left(\frac{\cos\theta_R + \dfrac{v}{c}}
        {1 + \dfrac{v}{c}\cos\theta_R}\right)^2\\
    &= \frac{1 + 2\dfrac{v}{c}\cos\theta_R
        + \left(\dfrac{v}{c}\right)^2\cos^2\theta_R - \cos^2\theta_R
        - 2\dfrac{v}{c}\cos\theta_R - \left(\dfrac{v}{c}\right)^2}
        {\left(1 + \dfrac{v}{c}\cos\theta_R\right)^2}\\
    &= \frac{\left(1 - \left(\dfrac{v}{c}\right)^2\right)
        \left(1 - \cos^2\theta_R\right)}
        {\left(1 + \dfrac{v}{c}\cos\theta_R\right)^2}
     = \frac{\sin^2\theta_R}
        {\gamma^2\left(1 + \dfrac{v}{c}\cos\theta_R\right)^2}\\
    \sin\theta_L &= \frac{\sin\theta_R}
        {\gamma\left(1 + \dfrac{v}{c}\cos\theta_R\right)}
  \end{align*}
  \hfill(3.0 pt)
  \begin{align*}
    \text{Also, } \Delta\theta_L \approx \sin(\Delta\theta_L)
      &= \frac{\sin(\Delta\theta_R)}
         {\gamma\left(1 + \dfrac{v}{c}\cos\theta_R\right)}\\
    \Delta\theta_L &\approx \frac{\Delta\theta_R}
      {\gamma\left(1 + \dfrac{v}{c}\cos\theta_R\right)}
  \end{align*}
  \hfill(2.0 pt)
  \begin{align*}
    \frac{\Delta\Omega_L}{\Delta\Omega_R}
      &= \frac{\sin\theta_L(\Delta\theta_L)(\Delta\phi)}
              {\sin\theta_R(\Delta\theta_R)(\Delta\phi)}\\
      &= \frac{\sin\theta_R(\Delta\theta_R)}
              {\gamma^2\left(1 + \dfrac{v}{c}\cos\theta_R\right)^2}
         \cdot\frac{1}{\sin\theta_R(\Delta\theta_R)}\\
    \therefore\ \Delta\Omega_L
      &= \frac{\Delta\Omega_R}
              {\gamma^2\left(1 + \dfrac{v}{c}\cos\theta_R\right)^2}
  \end{align*}
  \hfill(2.0 pt)
\item[(e)] Let $N(\theta_L)\Delta\Omega_L\Delta T$ be the number of photons
  arriving in the vicinity of $P$ within the element of solid angle
  $\Delta\Omega_L$ and time interval $T, T + \Delta T$ (it is important that
  $\Delta T \neq \Delta t$ where $\Delta t$ is the emission time of the same
  photons (in lab frame)). These being photons emitted by K into the solid
  angle $\Delta\Omega$ in the time interval $t_{0R}, t_{0R} + \Delta t_{0R}$,
  so that:
  \[ N(\theta_L)\Delta\Omega_L\Delta T = N\Delta\Omega_R\Delta t_{0R} \]
  \hfill(2.0 pt)

  From an earlier part:
  \begin{align*}
    cT &= ct + R - r\sin\theta_L = ct + R - r\sin(\omega t)\\
    % sic: the official solution writes sin(theta_L) here; from part (c) the
    % angle is theta = omega*t, not theta_L.
    \therefore\ c\Delta T &= c\Delta t
      - r\left[\sin(\omega t)\cos(\omega\Delta t)
      + \sin(\omega\Delta t)\cos(\omega t) - \sin(\omega t)\right]\\
    c\Delta T &= c\Delta t - r\omega\cos(\omega t)\Delta t
  \end{align*}
  \hfill(2.0 pt)
  \begin{align*}
    \frac{\Delta T}{\Delta t_{0R}}
      &= \frac{\Delta t}{\Delta t_{0R}}\frac{\Delta T}{\Delta t}
       = \gamma\left(1 - \frac{v}{c}\cos(\omega t)\right)\\
    \frac{\Delta t_{0R}}{\Delta T}
      &= \frac{1}{\gamma\left(1 - \dfrac{v}{c}\cos(\omega t)\right)}
  \end{align*}
  \hfill(3.0 pt)
  \begin{align*}
    \frac{N(t_L)}{N_R}
      &= \frac{\Delta\Omega_R\Delta t_{0R}}{\Delta\Omega_L\Delta T}\\
      &= \frac{\gamma^2\left(1 + \dfrac{v}{c}\cos\theta_R\right)^2}
              {\gamma\left(1 - \dfrac{v}{c}\cos(\omega t)\right)}\\
      &= \frac{\gamma\left(1 + \left(\dfrac{v}{c}\right)
          \dfrac{\cos\theta_L - \dfrac{v}{c}}
                {1 - \dfrac{v}{c}\cos\theta_L}\right)^2}
              {\left(1 - \dfrac{v}{c}\cos\theta_L\right)}\\
      &= \frac{\gamma}
              {\gamma^4\left(1 - \dfrac{v}{c}\cos\theta_L\right)^3}\\
    N(t_L) &= \frac{N_R}
      {\gamma^3\left(1 - \dfrac{v}{c}\cos\left(t_L - R/c\right)\right)^3}
    % sic: the official solution omits the factor omega inside the cosine;
    % it should read cos[omega(t_L - R/c)].
  \end{align*}
  \hfill(3.0 pt)

  In terms of energy fluxes
  \begin{align*}
    F_0 &= \frac{hf_R N_R}{R^2} = \frac{L}{4\pi R^2}\\
    F(t_L) &= \frac{hf_L N(t_L)}{R^2}
  \end{align*}
  \hfill(2.0 pt)
  \begin{align*}
    \frac{F(t_L)}{F_0} &= \frac{hf_L N(t_L)}{hf_R N_R}\\
      &= \frac{1}{\gamma^4\left(1 - \dfrac{v}{c}
         \cos\left[\omega\left(t - \dfrac{R}{c}\right)\right]\right)^4}\\
    F(t_L) &= \frac{L}{4\pi R^2\gamma^4\left(1 - \dfrac{v}{c}
      \cos\left[\omega\left(t - \dfrac{R}{c}\right)\right]\right)^4}
  \end{align*}
  \hfill(3.0 pt)

  The radiation is strongly beamed in the direction of motion of the source
  so that a remote observer in or near the orbital plane of the source sees
  strongly pulsed radiation.
\item[(f)] The amplification (for a given $v$) is highest when
  $\cos\theta_L = 1$. \hfill(1.0 pt)
  \begin{align*}
    F(t_L) &= \frac{F_0}{\left[\gamma\left(1 - \dfrac{v}{c}\right)\right]^4}\\
    A_{\max} &= \frac{1}{\left[0.0975 \times 0.05\right]^4}\\
    A_{\max} &\approx 1.8 \times 10^9
    % sic: the official solution substitutes 0.0975 = 1/gamma^2 where the
    % formula above calls for gamma = 3.2; evaluating [gamma(1-v/c)]^{-4}
    % directly gives ~1.5e3, not 1.8e9 (and similarly for A_min below).
  \end{align*}
  \hfill(1.0 pt)
  \begin{align*}
    \text{Similarly, } A_{\min}
      &= \frac{1}{\left[0.0975 \times 1.95\right]^4}\\
    A_{\min} &\approx 770
  \end{align*}
  \hfill(1.0 pt)
\end{description}

\solhead{10.15}{White Dwarf}
% source id: 2024_TH_T4
\begin{description}
\item[(a)] Simplifying the expression given in the problem statement:
  \[ \frac{\mathrm{d}P}{\mathrm{d}r} = -\frac{GM(r)\rho}{r^2} \]
  At $r = R$, it is possible to use the approximation provided and the
  expression $\rho = 3M/(4\pi R^3)$:
  \begin{align*}
    -\frac{P_c}{R} &= -\frac{GM\{3M/(4\pi R^3)\}}{R^2}\\
    P_c &= \frac{3GM^2}{4\pi R^4}
  \end{align*}
  \hfill(1.5 pt)

  Furthermore, electron density $n_e$ can be related to mass density as:
  \[ \rho = \mu_e m_p n_e \]
  \hfill(1.5 pt)

  Using now the equation of state:
  \begin{align*}
    P_c &= \left(\frac{3}{8\pi}\right)^{2/3}\frac{h^2}{5m_e}
      \left(\frac{\rho}{\mu_e m_p}\right)^{5/3}\\
    &= \left(\frac{3}{8\pi}\right)^{2/3}\frac{h^2}{5m_e}
      \left(\frac{3M}{4\pi R^3\mu_e m_p}\right)^{5/3}\\
    \frac{3GM^2}{4\pi R^4}
    &= \left(\frac{3}{8\pi}\right)^{2/3}\frac{h^2}{5m_e}
      \left(\frac{3}{4\pi\mu_e m_p}\right)^{5/3}\frac{M^{5/3}}{R^5}\\
    R &= \left(\frac{4\pi}{3}\right)\left(\frac{3}{8\pi}\right)^{2/3}
      \frac{h^2}{5Gm_e}\left(\frac{3}{4\pi\mu_e m_p}\right)^{5/3} M^{-1/3}
  \end{align*}
  Therefore:
  \[ \boxed{a = \left(\frac{4\pi}{3}\right)
     \left(\frac{3}{8\pi}\right)^{2/3}\frac{h^2}{5Gm_e}
     \left(\frac{3}{4\pi\mu_e m_p}\right)^{5/3}} \]
  \hfill(1.5 pt)
  \[ \boxed{b = -\frac{1}{3}} \]
  \hfill(1.5 pt)
\item[(b)] There are 2 nucleons (1 proton and 1 neutron) per unit electron
  for carbon, so that $\mu_e = 2$. \hfill(1.0 pt)

  Using the expression for the radius found on the previous item:
  \begin{align*}
    R &\approx \left(\frac{4\pi}{3}\right)
      \left(\frac{3}{8\pi}\right)^{2/3}\frac{1}{5}
      \left(\frac{3}{4\pi}\right)^{5/3}\frac{h^2}{Gm_e}\,
      \frac{1}{(\mu_e m_p)^{5/3}}\,\frac{1}{M^{1/3}}\\
    &\approx 1.866\cdot10^{-2}\,\frac{h^2}{Gm_e}\,
      \frac{1}{(\mu_e m_p)^{5/3}}\,\frac{1}{M^{1/3}}\\
    &\approx \frac{1.866\times10^{-2}\times
      \left(6.626\times10^{-34}\right)^2}
      {6.67\times10^{-11}\times 9.11\times10^{-31}\times
      \left(2\times1.67\times10^{-27}\right)^{5/3}\times
      \left(1.988\times10^{30}\right)^{1/3}}
  \end{align*}
  \[ \boxed{R \approx 1.44\times10^{6}\ \mathrm{m}} \]
  \hfill(3.0 pt)
\end{description}

\solhead{10.16}{Physics of Accretion}
% source id: 2024_TH_T9
\begin{description}
\item[(a)] For a Keplerian orbit:
  \[ \Delta m\frac{v_k^2}{r} = \frac{GM\Delta m}{r^2}
     \ \Rightarrow\ v_k^2 = \frac{GM}{r} \]
  \hfill(2.0 pt)

  Total energy:
  \[ E = E_{kin} + E_{pot}
     \ \Rightarrow\ E(r) = \frac{1}{2}\Delta m v_k^2 - \frac{GM\Delta m}{r}
     \ \Rightarrow\ E(r) = -\frac{1}{2}\frac{GM\Delta m}{r} \]
  \hfill(2.0 pt)
  \[ E(R_{max}) - E(r)
     = \frac{1}{2}GM\Delta m\left(\frac{1}{r} - \frac{1}{R_{max}}\right)
     = \frac{GM\Delta m}{2r}\left(1 - \frac{r}{R_{max}}\right)
     \approx \frac{GM\Delta m}{2r} \]
  \hfill(1.0 pt)

  Therefore:
  \[ \boxed{\frac{\Delta E}{\Delta m} \approx \frac{GM}{2r}} \]
  \hfill(1.0 pt)
\item[(b)] Variation in energy:
  \[ \Delta E_{Tot} = E(R_{max}) - E(R_{min}) \]
  Luminosity:
  \[ L_{Tot} = \frac{\Delta E_{Tot}}{\Delta t}
     = \frac{\Delta E_{Tot}}{\Delta m}\times\dot{M} \]
  \hfill(3.0 pt)
  \[ \boxed{L_{Tot} = \frac{GM\dot{M}}{2R_{min}}} \]
  \hfill(2.0 pt)

  Where the result from (a) was used, allied with $R_{min} \ll R_{max}$.
\item[(c)] Variation in the energy of the tiny mass:
  \[ \Delta E = E(r + \Delta r) - E(r)
     = \frac{1}{2}GM\Delta m\left(\frac{1}{r} - \frac{1}{r + \Delta r}\right)
     \approx \frac{1}{2}GM\Delta m
       \left(\frac{\Delta r}{r^2(1 + \Delta r/r)}\right) \]
  \[ \Rightarrow\ \frac{\Delta E}{\Delta m}
     \approx \frac{1}{2}GM\Delta r\,\frac{1}{r^2} \]
  \hfill(5.0 pt)

  Multiplying this expression by the rate of mass variation per time, it is
  possible to obtain an expression for $\frac{\Delta E}{\Delta t\Delta r}$:
  \[ \frac{\Delta E}{\Delta t} = \frac{\Delta E}{\Delta m}\times\dot{M}
     = \frac{GM\dot{M}\Delta r}{2r^2} \]
  \[ \Rightarrow\ \boxed{\frac{\Delta E}{\Delta t\Delta r}
     = \frac{GM\dot{M}}{2r^2}} \]
  \hfill(3.0 pt)
\item[(d)] Surface area of the ring:
  \[ A = 2\left[\pi(r + \Delta r)^2 - \pi r^2\right]
     = 2(\pi r^2 + 2\pi r\Delta r + \pi\Delta r^2 - \pi r^2) \]
  \hfill(4.0 pt)
  \[ \Rightarrow\ A \approx 4\pi r\Delta r \]
  Using Stefan-Boltzmann's law:
  \[ A\sigma T^4 = L \]
  \hfill(2.0 pt)
  \[ 4\pi r\Delta r\sigma T^4 = \frac{\Delta E}{\Delta t}
     \ \Rightarrow\ 4\pi r\sigma T^4 = \frac{\Delta E}{\Delta r\Delta t} \]
  \hfill(2.0 pt)
  \begin{align*}
    4\pi r\sigma T^4 &= \frac{GM\dot{M}}{2r^2}\\
    \Rightarrow\ &\boxed{T = \left(\frac{GM\dot{M}}
      {8\pi\sigma r^3}\right)^{1/4}}
  \end{align*}
  \hfill(2.0 pt)
\item[(e)] Internal radius of the ring:
  \[ R_{sch} = \frac{2GM}{c^2}
     \ \Rightarrow\ R_{min} = 3R_{sch} = \frac{6GM}{c^2} \]
  \[ R_{min} = \frac{6GM}{c^2}
     = \frac{6\times 6.674\times10^{-11}\times 3\times 1.988\times10^{30}}
            {(2.998\times10^8)^2}
     = 2.7\times10^4\ \mathrm{m} \]
  Total luminosity:
  \[ L_{Tot} = \frac{GM\dot{M}}{2R_{min}} = \frac{\dot{M}c^2}{12} \]
  \[ \Rightarrow\ L_{Tot}
     = \frac{10^{-9}\times 1.988\times10^{30}\times(2.998\times10^8)^2}
            {12\times 365.2422\times 24\times 60\times 60} \]
  \[ \Rightarrow\ \boxed{L_{Tot} = 5\times10^{29}\ \mathrm{W}} \]
  \hfill(1.5 pt)

  Surface temperature:
  \[ T = \left(\frac{GM\dot{M}}{8\pi\sigma R_{min}^3}\right)^{1/4}
     = \left(\frac{6.674\times10^{-11}\times 3\times10^{-9}\times
       (1.988\times10^{30})^2}
       {8\pi\times 5.670\times10^{-8}\times(2.7\times10^4)^3\times
       365.2422\times 24\times 60\times 60}\right)^{1/4} \]
  \[ T = 5.5\times10^6\ \mathrm{K} \]
  Using Wien's law:
  \[ \lambda = \frac{b}{T} = \frac{2.898\times10^{-3}}{5.5\times10^6} \]
  \[ \Rightarrow\ \boxed{\lambda = 5\times10^{-10}\ \mathrm{m}} \]
  \hfill(1.5 pt)
\item[(f)] Using Wien's law:
  \[ \lambda = \frac{b}{T}
     \ \Rightarrow\ T = \frac{b}{\lambda}
     = \frac{2.898\times10^{-3}}{6\times10^{-8}} \]
  \[ \Rightarrow\ T = 4.8\times10^4\ \mathrm{K} \]
  \hfill(1.0 pt)

  Considering this temperature to be that of the innermost of the accretion
  disk, and using $R_{\min} = 3R_{\mathrm{sch}}$:
  \[ T = \left(\frac{GM\dot{M}}{8\pi\sigma R_{min}^3}\right)^{1/4}
     \ \Rightarrow\ T^4
     = \left(\frac{GM\dot{M}}{8\pi\sigma(3R_{\mathrm{sch}})^3}\right)
     = \left(\frac{GM\dot{M}}{8\pi\sigma(6GM/c^2)^3}\right) \]
  \[ \Rightarrow\ M = \left(\frac{\dot{M}}{8\pi\sigma 6^3}\right)^{1/2}
     \frac{c^3}{GT^2} \]
  \[ \Rightarrow\ M = \frac{1}{(8\pi)^{1/2}}\,\frac{1}{6^{3/2}}
     \left\{\frac{1.988\times10^{30}/(3.15\times10^7)}
       {5.670\times10^{-8}}\right\}^{1/2}
     \frac{(2.998\times10^8)^3}
       {6.674\times10^{-11}\times(4.8\times10^4)^2} \]
  \[ \boxed{M \approx 2.5\times10^{39}\ \mathrm{kg}
     \approx 1.3\times10^9\ M_\odot} \]
  \hfill(2.0 pt)
\end{description}

\solhead{10.17}{Neutron Star Binary}
% source id: 2025_TH_T07
\begin{description}
\item[(T07.1)] The accretion radius can be estimated by equating the
  gravitational potential energy of the stellar wind to its kinetic energy
  \begin{align*}
    \frac{1}{2}\rho_{\mathrm{w}} v_{\mathrm{rel}}^2
      &= \frac{G M_{\mathrm{NS}} \rho_{\mathrm{w}}}{R_{\mathrm{acc}}}
      \tag*{(0.5)}\\
    \implies R_{\mathrm{acc}} &= \frac{2 G M_{\mathrm{NS}}}{v_{\mathrm{rel}}^2}
      \tag*{(0.5)}
  \end{align*}
  where $\rho_{\mathrm{w}}$ is the wind density and $v_{\mathrm{rel}}$ is the
  relative velocity of the wind.

  For the case considered above,
  \[
    v_{\mathrm{rel}}^2 = v_{\mathrm{w}}^2 + v_{\mathrm{orb}}^2
  \]
  Here, $v_{\mathrm{w}} \gg v_{\mathrm{orb}}$, therefore,
  $v_{\mathrm{rel}}^2 \approx v_{\mathrm{w}}^2
  = 9 \times 10^{12}\,\mathrm{m^2\,s^{-2}}$. \hfill(1.0)

  which gives,
  \[
    R_{\mathrm{acc}} = \frac{2 \times 6.674 \times 10^{-11} \times 2.0 \times
      1.988 \times 10^{30}}{9 \times 10^{12}}\,\mathrm{m}
  \]
  \[
    R_{\mathrm{acc}} = 5.9 \times 10^{4}\,\mathrm{km}
  \]
  \hfill(1.0)

\item[(T07.2)] In the case of spherically symmetric mass loss from the
  companion and $v_{\mathrm{w}} \gg v_{\mathrm{orb}}$, the companion mass loss
  is
  \[
    \dot{M}_{\mathrm{OB}} = 4\pi a^2 \rho_{\mathrm{w}} v_{\mathrm{w}}
  \]
  \hfill(1.0)

  where $a = 0.5\,\mathrm{au} = 7.480 \times 10^{10}$\,m is the orbital
  separation. Mass accretion rate on the compact object is
  \[
    \dot{M}_{\mathrm{acc}} \approx \pi R_{\mathrm{acc}}^2\,\rho_{\mathrm{w}}\,
      v_{\mathrm{rel}}
  \]
  \hfill(1.0)

  Also here, since $v_{\mathrm{w}} \gg v_{\mathrm{orb}}$,
  $v_{\mathrm{rel}} \approx v_{\mathrm{w}}$. This implies,
  \begin{align*}
    \dot{M}_{\mathrm{acc}}
      &\approx \left(\frac{\dot{M}_{\mathrm{OB}}}{4}\right)
        \left(\frac{R_{\mathrm{acc}}}{a}\right)^2
        \approx 9.74 \times 10^{11}\,\mathrm{kg\,s^{-1}}\\
    \implies \dot{M}_{\mathrm{acc}}
      &\approx 1.6 \times 10^{-11}\,\mathrm{M_\odot\,yr^{-1}}
      \tag*{(1.0)}
  \end{align*}

\item[(T07.3)] The orbital speed of the NS in the frame corresponding to the
  center of OB star (assuming the star is stationary) is given by
  \[
    v_{\mathrm{orb}}^2 = \frac{G M_{\mathrm{OB}}}{a}
  \]
  \hfill(1.0)

  Using vector algebra we get
  $\tan\beta = v_{\mathrm{orb}}/v_{\mathrm{w}}$. \hfill(1.0)

  Now, the ratio of the mass accretion rate to that of mass loss rate from the
  star
  \begin{align*}
    \frac{\dot{M}_{\mathrm{acc}}}{\dot{M}_{\mathrm{OB}}}
      &\sim \frac{\pi R_{\mathrm{acc}}^2 \rho_{\mathrm{w}} v_{\mathrm{rel}}}
        {4\pi a^2 \rho_{\mathrm{w}} v_{\mathrm{w}}}\\
      &= \frac{1}{4}\left(\frac{R_{\mathrm{acc}}}{a}\right)^2
        \sqrt{\frac{v_{\mathrm{orb}}^2 + v_{\mathrm{w}}^2}{v_{\mathrm{w}}^2}}
       = \frac{1}{4}\left(\frac{R_{\mathrm{acc}}}{a}\right)^2
        \sqrt{1 + \tan^2\beta}
      \tag*{(1.0)}
  \end{align*}
  To find the expression of the ratio of the two radii, $R_{\mathrm{acc}}$ and
  orbital separation $a$,
  \begin{align*}
    \frac{R_{\mathrm{acc}}}{a}
      &= \frac{2 G M_{\mathrm{NS}}}{v_{\mathrm{rel}}^2}\frac{1}{a}
       = \frac{2 G M_{\mathrm{OB}}\,q}{v_{\mathrm{rel}}^2}\frac{1}{a}
       = \frac{2 v_{\mathrm{orb}}^2}{v_{\mathrm{rel}}^2}\,q
      \tag*{(1.0)}\\
      &= \left(\frac{2\tan^2\beta}{1 + \tan^2\beta}\right) q
      \tag*{(1.0)}
  \end{align*}
  On substituting and simplifying, we can express,
  \[
    \implies f(\tan\beta, q)
      = q^2 \left(\frac{\tan^4\beta}{(1 + \tan^2\beta)^{3/2}}\right)
  \]
  \hfill(1.0)

\item[(T07.4a)] The magnetic pressure $P_{\mathrm{mag}} = B^2/2\mu_0$
  increases rapidly towards the NS magnetic poles.
  \[
    \implies P_{\mathrm{eq,\,mag}}
      = \frac{B_0^2 R_{\mathrm{NS}}^6}{2\mu_0 r^6}
  \]
  \hfill(1.0)

\item[(T07.4b)] The critical magnetic field $B_{0,\mathrm{c}}$ can be obtained
  by equating pressure due to incoming wind with the $P_{\mathrm{eq,\,mag}}$ at
  $r = R_{\mathrm{acc}}$. \hfill(1.0)

  The pressure due to incoming wind at $r = R_{\mathrm{acc}}$ is given by
  \[
    \rho_{\mathrm{w}} v_{\mathrm{rel}}^2 \sim \rho_{\mathrm{w}} v_{\mathrm{w}}^2
  \]
  \hfill(1.0)

  From part (T07.2), one can obtain the expression for $\rho_{\mathrm{w}}$ as
  follows:
  \[
    \rho_{\mathrm{w}}
      = \frac{\dot{M}_{\mathrm{acc}}}{\pi R_{\mathrm{acc}}^2 v_{\mathrm{w}}}
      = \frac{\dot{M}_{\mathrm{OB}}}{4\pi a^2 v_{\mathrm{w}}}
  \]
  \hfill(1.0)

  Now equating the pressures due to wind and magnetic field, we have:
  \[
    \frac{B_{0,\mathrm{c}}^2 R_{\mathrm{NS}}^6}{2\mu_0 R_{\mathrm{acc}}^6}
      = \frac{\dot{M}_{\mathrm{acc}}}{\pi R_{\mathrm{acc}}^2}
        \left(\frac{2 G M_{\mathrm{NS}}}{R_{\mathrm{acc}}}\right)^{1/2}
  \]
  \hfill(1.0)

  which gives the value of the critical magnetic field as
  \[
    \implies B_{0,\mathrm{c}}
      = \left(\frac{2\mu_0 \dot{M}_{\mathrm{acc}}
          \sqrt{2 G M_{\mathrm{NS}}}}{\pi R_{\mathrm{NS}}^6}\right)^{1/2}
        R_{\mathrm{acc}}^{7/4}
  \]
  \hfill(1.0)

  Now putting the values we get
  \[
    \implies B_{0,\mathrm{c}} \approx 4.0 \times 10^{9}\,\mathrm{T}
  \]
  \hfill(2.0)
\end{description}

\solhead{10.18}{Shadow of a black hole}
% source id: 2025_TH_T08
\begin{description}
\item[(T08.1)] Circular trajectories have $r = $ constant. Hence, the radial
  velocity $v_{\mathrm{r}} = 0$ \hfill(1.0)

  and rate of change of radial component of velocity
  $\dfrac{dv_{\mathrm{r}}}{dt} = 0$. \hfill(1.0)

  From the third equation,
  \[
    \frac{dv_{\mathrm{r}}}{dt} - \frac{L^2}{r^3} + \frac{3 G M L^2}{c^2 r^4}
    = 0,
  \]
  we get radius $r_{\mathrm{ph}}$ of circular orbits to be
  $r_{\mathrm{ph}} = 3GM/c^2$. \hfill(1.0)

  Then from the first equation, the impact parameter is found to be
  \[
    b_{\mathrm{ph}} = L/\sqrt{2E} = 3\sqrt{3}\,GM/c^2.
  \]
  \hfill(1.0)

\item[(T08.2)] Time taken to complete one circular orbit is
  \begin{align*}
    T_{\mathrm{ph}} &= 2\pi r_{\mathrm{ph}}/c = 6\pi G M/c^3
      \tag*{(1.0)}\\
      &= 6.0 \times 10^{5}\,\mathrm{s}
      \tag*{(1.0)}
  \end{align*}

\item[(T08.3a)] $r_\alpha$ corresponds to the radius where
  $V_{\mathrm{eff}} = \dfrac{L^2}{2r^2}\left(1 - \dfrac{2GM}{c^2 r}\right)$ is
  equal to zero. Hence $r_\alpha = 2GM/c^2 = R_{\mathrm{SC}}$. \hfill(1.0)

  Solving for $\dfrac{dV_{\mathrm{eff}}}{dr} = 0$ gives
  $r_\beta = 3GM/c^2 = r_{\mathrm{ph}}$. \hfill(1.0)

\item[(T08.3b)] The radius of the position of peak of $V_{\mathrm{eff}}(r)$ is
  the smallest possible turning point radius $r_{\mathrm{t}}$ for the light
  trajectory such that it can escape to infinity. \hfill(1.0)

  Hence, the least smallest possible turning point radius $r_{\mathrm{t}}$ is
  same as the radius $r_{\mathrm{ph}}$ of the circular orbits.
  $r_{\mathrm{t}} = 3GM/c^2$ \hfill(1.0)

  Since $r_{\mathrm{t}}$ is equal to $r_{\mathrm{ph}}$, the corresponding
  impact parameter is same as calculated in part one
  $b_{\mathrm{min}} = 3\sqrt{3}\,GM/c^2$. \hfill(1.0)

\item[(T08.4)] Since it is given that $r_{\mathrm{app}} \approx b$, the
  apparent radii $r_{\mathrm{inner}}$ and $r_{\mathrm{outer}}$ are related to
  the respective impact parameters $b_{\mathrm{inner}}$ and
  $b_{\mathrm{outer}}$ of the corresponding light rays.

  To figure out the $b_{\mathrm{inner}}$ and $b_{\mathrm{outer}}$, one needs to
  choose the correct corresponding $r_{\mathrm{actual}}$ in each case.

  The outer ring radius $r_{\mathrm{outer}}$ would correspond to the light ray
  which is emitted from the outer-edge of the accretion disk. Hence the
  $r_{\mathrm{outer}}$ is determined by the formula given
  $r_{\mathrm{outer}} \approx b_{\mathrm{outer}} \approx a_{\mathrm{outer}}
  \left(1 + R_{\mathrm{SC}}/a_{\mathrm{outer}}\right)^{1/2}$. \hfill(1.0)

  Substituting $a_{\mathrm{outer}} = 10 R_{\mathrm{SC}}$ we get
  $r_{\mathrm{outer}} = 1.3 \times 10^{3}$\,au. \hfill(1.0)

  For the inner ring radius $r_{\mathrm{inner}}$, the light ray corresponding
  to the smallest turning point $r_{\mathrm{t}}$ is the light ray which goes
  closest to the BH and returns back to infinity.

  Thus one should choose the impact parameter $b_{\mathrm{inner}}$ to be that
  of the circular orbit. Then,
  $r_{\mathrm{inner}} \approx b_{\mathrm{inner}} = b_{\mathrm{ph}}
  = 3\sqrt{3}\,GM/c^2$. \hfill(2.0)

  Hence $r_{\mathrm{inner}} = 3.3 \times 10^{2}$\,au. \hfill(1.0)

\item[(T08.5a)] The source is a point source.

  For the image, one has to look at all the possible trajectories light that
  can reach the observer starting from point Z in the same plane as the
  observer, the black hole and point Z.

  The shortest path is the direct path, without orbiting the black hole,
  connecting Z and the observer. This will be also be bent due to
  gravity. % sic: "will be also be" in source
  The
  next shortest path could be the one which goes half around the black hole in
  the anti-clockwise sense in the picture. Another one is the one with one and
  half revolution around the black hole in the anti-clockwise sense in the
  picture. And so on, paths with additional one more revolution around the
  black hole than the previous one.

  The above is also true for trajectories in clockwise sense of rotation
  around the black hole. Therefore, number of possible paths for light to
  travel from Z to the observer is (D) Greater than two. \hfill(2.0)

\item[(T08.5b)] One has to consider the time difference in the arrival times
  due to the different path lengths. From Part (T08.2), it takes around
  $6.0 \times 10^{5}$\,sec which is 166.7\,h for light to circularly orbit the
  supermassive black hole of mass $M = 6.5 \times 10^{9}\,\mathrm{M_\odot}$.
  Hence, the difference in the arrival times of light along these trajectories
  would be at least around half of 166.7\,h approximately, that is, around
  83.3\,h. Each of these signals would last for only the duration of the
  blast, i.e., 5\,s.

  Since the long exposure picture is only for 60\,s, either one or no image
  will be captured, depending on whether the exposure time encompasses the
  short window of the arrival of one of the signals and 5\,s thereafter.
  Therefore, the number of images in a picture would be (A) At most one.
  \hfill(2.0)
\end{description}

\solhead{10.19}{Hawking Radiation from Black Holes}
% source id: 2025_TH_T12
\begin{description}
\item[(T12.1a)] Total energy of any particle is
  $E = E_{\mathrm{grav}} + mc^2$, so we have
  \[
    E_1 = -\frac{G M_{\mathrm{bh}} m}{R_{\mathrm{SC}}} + mc^2, \qquad
    E_2 = -\frac{G M_{\mathrm{bh}} m}{\kappa R_{\mathrm{SC}}} + mc^2.
  \]
  \hfill(1.0)

  We require $E_1 + E_2 = 0$. Hence,
  \[
    2mc^2 - \frac{G M_{\mathrm{bh}} m}{R_{\mathrm{SC}}}
    \left(1 + \frac{1}{\kappa}\right) = 0
  \]
  \hfill(1.0)

  Then
  \[
    \frac{1}{\kappa} = \frac{2 c^2 R_{\mathrm{SC}}}{G M_{\mathrm{bh}}} - 1 = 3
  \]
  where we have used
  $R_{\mathrm{SC}} = \dfrac{2 G M_{\mathrm{bh}}}{c^2}$. \hfill(1.0)

  This gives
  \[
    \kappa = \frac{1}{3}.
  \]
  \hfill(1.0)

\item[(T12.1b)] Using Wien's law,
  $\lambda_{\mathrm{bb}} = b/T_{\mathrm{bh}}$ where
  $b = 2.898 \times 10^{-3}$\,m\,K. \hfill(1.0)

  Hence,
  \[
    T_{\mathrm{bh}} = \frac{b}{16 R_{\mathrm{SC}}}
    \implies T_{\mathrm{bh}} = \frac{b c^2}{32 G M_{\mathrm{bh}}}
  \]
  \hfill(1.0)

  Since $R_{\mathrm{SC}}$ for a solar mass black hole is 2.952\,km, for a
  black hole with mass $= 10\,M_\odot$, we get
  $R_{\mathrm{SC,\,10\odot}} \simeq 29.52$\,km. \hfill(1.0)

  The temperature of this BH will be
  \[
    T_{\mathrm{bh,\,10\odot}}
    = \frac{2.898 \times 10^{-3}\,\mathrm{m\,K}}{16 \times
      R_{\mathrm{SC,\,10\odot}}}
  \]
  \[
    T_{\mathrm{bh,\,10\odot}} = 6.1 \times 10^{-9}\,\mathrm{K}.
  \]
  \hfill(1.0)

\item[(T12.1c)] Power emitted by an evaporating black hole
  $= \sigma T_{\mathrm{bh}}^4 4\pi R_{\mathrm{SC}}^2$. \hfill(1.0)

  This gives
  \[
    \frac{dM_{\mathrm{bh}}}{dt}
    = -\frac{1}{c^2}\sigma T_{\mathrm{bh}}^4 4\pi R_{\mathrm{SC}}^2.
  \]
  \hfill(1.0)

  With $T_{\mathrm{bh}} = b/(16 R_{\mathrm{SC}})$,
  \[
    \frac{dM_{\mathrm{bh}}}{dt}
    = -\frac{\pi b^4 \sigma c^2}{16^4 G^2}\,\frac{1}{M_{\mathrm{bh}}^2}
  \]
  \hfill(1.0)

  This gives
  $\dfrac{dM_{\mathrm{bh}}}{M_{\mathrm{bh}}^{2}}
  = -\dfrac{\pi b^4 \sigma c^2}{16^4 G^2}\,dt$,
  % sic: source writes dM_bh/M_bh^2, though the integral that follows is of
  % M_bh^2 dM_bh
  and hence
  $\dfrac{M_{\mathrm{bh}}^3}{3}
  = -\dfrac{\pi b^4 \sigma c^2}{16^4 G^2}\,t + C$. \hfill(1.0)
  Since $M_{\mathrm{bh}} = M_0$ at $t = 0$, we get
  \[
    M_{\mathrm{bh}}^3 = M_0^3 - 3\frac{\pi b^4 \sigma c^2}{16^4 G^2} t
  \]
  \hfill(0.5)
  \[
    M_{\mathrm{bh}}(t)
    = \left(M_0^3 - 3\frac{\pi b^4 \sigma c^2}{16^4 G^2} t\right)^{1/3}
  \]
  \hfill(1.0)

  The sketch of $M_{\mathrm{bh}}(t)$ is a convex curve starting at $M_0$ on
  the $M_{\mathrm{bh}}(t)$-axis, with negative initial slope and large
  negative final slope, reaching zero at
  $t = \dfrac{16^4 G^2 M_0^3}{3\pi b^4 \sigma c^2}$. \hfill(2.5)
  \ioaafig{2025_TH_T12-s1.pdf}

\item[(T12.1d)] We know that
  $M(t) = \left(M_0^3 - 3\dfrac{\pi b^4 \sigma c^2}{16^4 G^2}
  t\right)^{1/3}$.

  Clearly, $M(t) = 0$ when
  $t = \dfrac{M_0^3}{3}\,\dfrac{16^4 G^2}{\pi b^4 \sigma c^2}$. Thus,
  \[
    \tau_{\mathrm{bh}} = \frac{M_0^3}{3}\,
      \frac{16^4 G^2}{\pi b^4 \sigma c^2}
  \]
  \hfill(1.0)

  With $A \equiv \dfrac{\pi b^4 \sigma c^2}{16^4 G^2}
  \simeq 3.868 \times 10^{15}\,\mathrm{kg^3\,s^{-1}}$, this gives \hfill(0.5)
  \[
    \tau_{\mathrm{bh}} \approx \frac{M_0^3}{3A}
    = 8.617 \times 10^{-17} \left(\frac{M_0}{\mathrm{kg}}\right)^3
    \mathrm{s}.
  \]
  \hfill(0.5)

  For $M_0 = 10\,\mathrm{M_\odot}$,
  \[
    \tau_{\mathrm{bh,\,10\odot}} \approx 6.8 \times 10^{77}\,\mathrm{s}
  \]
  \hfill(1.0)

\item[(T12.1e)] As the CMB will redshift away eventually, there will be no
  CMB radiation for the BH to absorb at large $t$. It will only decrease in
  mass and increase in temperature. So only the figures which show the
  temperature of the black hole increasing eventually ($t \to \infty$) are
  relevant. The derivatives at the origin (now) should distinguish the three
  cases.
  \begin{description}
  \item[(X)] If $T_{\mathrm{bh}}^{\mathrm{now}} >
    T_{\mathrm{cmb}}^{\mathrm{now}}$, then the BH is losing more energy now
    than it is gaining by CMB absorption, and its temperature should increase
    due to decreasing mass. $(dT/dt)_{\mathrm{now}} > 0$.
  \item[(Y)] If $T_{\mathrm{bh}}^{\mathrm{now}} =
    T_{\mathrm{cmb}}^{\mathrm{now}}$, the temperature will be in equilibrium
    for some time, i.e.\ $(dT/dt)_{\mathrm{now}} = 0$, before the CMB
    temperature redshifts to smaller values and $(dT/dt)_{\infty} > 0$.
  \item[(Z)] For $T_{\mathrm{bh}}^{\mathrm{now}} <
    T_{\mathrm{cmb}}^{\mathrm{now}}$, initially the BH will gain mass (and
    hence cool down) by absorbing from CMBR compared to what it loses by
    evaporation. So $(dT/dt)_{\mathrm{now}} < 0$. But eventually,
    $(dT/dt)_{\infty} > 0$.
  \end{description}
  Therefore: X $\to$ III, \quad Y $\to$ II, \quad Z $\to$ VI. \hfill(6.0)

\item[(T12.2a)] We want
  $\tau_{\mathrm{PBH}} = 14 \times 10^{9} \times 3.156 \times 10^{7}\,\mathrm{s}
  = 4.418 \times 10^{17}$\,s. \hfill(1.0)

  But we have found in (T12.1d) that
  $\tau_{\mathrm{bh}} = 8.617 \times 10^{-17}
  \left(\dfrac{M_0}{\mathrm{kg}}\right)^3$\,s. This gives
  \[
    M_{0,\mathrm{PBH}}^3
    = \left(\frac{\tau_{\mathrm{PBH}}}{8.617 \times 10^{-17}}\right)
      \mathrm{kg}^3
    \Rightarrow M_{0,\mathrm{PBH}} \simeq 1.7 \times 10^{11}\,\mathrm{kg}
  \]
  \hfill(1.0)

  Since $R_{\mathrm{SC}} \propto M_{\mathrm{bh}}$, and we know
  $R_{\mathrm{SC,\odot}} = 2952$\,m,
  \[
    R_{\mathrm{SC,\,PBH}} \simeq R_{\mathrm{SC,\odot}}
    \left(\frac{M_{0,\mathrm{PBH}}}{M_\odot}\right)
    \Rightarrow R_{\mathrm{SC,\,PBH}} = 2.6 \times 10^{-16}\,\mathrm{m}
  \]
  \hfill(1.0)

  Since $T_{\mathrm{bh}} \propto 1/M_{\mathrm{bh}}$ and we know
  $T_{\mathrm{bh,\,10\odot}}$ from (T12.1b),
  \[
    T_{\mathrm{PBH}} \simeq T_{\mathrm{bh,\,10\odot}}
    \left(\frac{10 M_\odot}{M_{0,\mathrm{PBH}}}\right)
    \Rightarrow T_{\mathrm{PBH}} = 7.1 \times 10^{11}\,\mathrm{K}
  \]
  \hfill(1.0)

\item[(T12.2b)] The size of horizon $r_{\mathrm{H}} = ct$. So the total mass
  contained within the horizon volume is
  \[
    M = \frac{1}{c^2}\rho \left(\frac{4\pi}{3}\right)
      \left(\frac{ct}{2}\right)^3.
  \]
  \hfill(2.0)

  With $a(t) \sim t^{1/2}$ and $a(t) \sim 1/T$, one gets
  $t \sim T^{-2}$. \hfill(2.0)

  With $\rho \sim T^4 \sim t^{-2}$, the mass of PBH formed at time $t$ would
  be
  \[
    M_{0,\mathrm{PBH}} \propto t.
  \]
  \hfill(1.0)

  The age of the Universe at which PBH of $10^{12}$\,kg forms is
  $t_{12} = 10^{-23}$\,s. Hence,
  \[
    t_{20} = t_{12} \times
    \left(\frac{10^{20}\,\mathrm{kg}}{10^{12}\,\mathrm{kg}}\right)
    \Rightarrow t_{20} = t_{12} \times 10^{8} = 10^{-15}\,\mathrm{s}
  \]
  \hfill(1.0)

\item[(T12.2c)] Lifetime of the PBH with
  $M_{0,\mathrm{PBH}} = 10^{10}$\,kg will be
  \[
    \tau_{\mathrm{PBH}} = 8.617 \times 10^{-17} \cdot (10^{10})^3\,\mathrm{s}
    = 8.617 \times 10^{13}\,\mathrm{s}.
  \]
  \hfill(0.5)

  As given in (T12.1b), $R_{\mathrm{SC,\odot}} = 2952$\,m. Using
  $R_{\mathrm{SC}} \propto M$, we get the peak wavelength of radiation
  emitted by this $10^{10}$\,kg PBH, at the time of emission, as
  \[
    \lambda_0 \simeq 16 R_{\mathrm{SC}}
    = 16 \cdot \frac{10^{10}\,\mathrm{kg}}{M_\odot} \cdot
      R_{\mathrm{SC,\odot}}
    = 2.4 \times 10^{-16}\,\mathrm{m}.
  \]
  \hfill(1.5)

  Redshifted wavelength received at earth at present time $t_0$ is obtained
  from
  \[
    \lambda_{\mathrm{earth}}
    = \lambda_0 \frac{a(t_0)}{a(\tau_{\mathrm{PBH}})}.
  \]
  \hfill(1.0)

  With $a(t) \sim t^{2/3}$, we get
  \[
    \lambda_{\mathrm{earth}}
    = \lambda_0 \left(\frac{t_{\mathrm{now}}}{\tau_{\mathrm{PBH}}}\right)^{2/3}
    = 2.4 \times 10^{-16}\,\mathrm{m}
      \left(\frac{4.418 \times 10^{17}\,\mathrm{s}}
        {8.617 \times 10^{13}\,\mathrm{s}}\right)^{2/3}
  \]
  \hfill(1.0)
  \[
    \lambda_{\mathrm{earth}} = 7.1 \times 10^{-14}\,\mathrm{m}
  \]
  \hfill(1.0)

\item[(T12.2d)] Using the expression for $T_{\mathrm{bh}}$ found in
  (T12.1b), the energy of a typical photon emitted at time $t$ for a PBH is
  \[
    E_{\mathrm{emitted}}(t) = k_{\mathrm{B}} T_{\mathrm{bh}}(t)
    = \frac{b c^2 k_{\mathrm{B}}}{32 G M_{\mathrm{bh}}(t)}.
  \]
  \hfill(1.0)

  From (T12.1c), we know
  \[
    M_0^3 = 3\frac{\pi b^4 \sigma c^2}{16^4 G^2} t + M_{\mathrm{bh}}(t)^3
    = 3At + \left(\frac{b c^2 k_{\mathrm{B}}}
      {32 G E_{\mathrm{emitted}}(t)}\right)^3.
  \]
  \hfill(1.0)

  The value of $E_{\mathrm{det}}$ can be different from
  $E_{\mathrm{emitted}}$ due to redshift. Hence, the smallest possible value
  of $E_{\mathrm{emitted}}(t)$ is
  $E_{\mathrm{det}} = 3.0 \times 10^{20}$\,eV. This gives the maximum value
  of the second term as
  \[
    \left(\frac{b c^2 k_{\mathrm{B}}}
      {32 G E_{\mathrm{emitted}}(t)}\right)^3_{\mathrm{max}}
    = 4.3 \times 10^{-5}\,\mathrm{kg}^3.
  \]
  \hfill(1.0)

  Given the value of $A = 3.868 \times 10^{15}\,\mathrm{kg^3\,s^{-1}}$, the
  second term can always be neglected for $t \gtrsim 10^{-19}$\,s. This is
  true in this problem. So $M_0 \approx (3At)^{1/3}$. \hfill(1.0)

  \textbf{The largest possible} $M_0$ would be obtained when $t$ is as large
  as possible, i.e., $t = t_{\mathrm{now}}$. This will correspond to a BH
  radiating in its last stages at present in the nearby region. \hfill(1.0)

  Then
  \[
    M_0^{\mathrm{max}} = (3At_{\mathrm{now}})^{1/3}
    = (5.13 \times 10^{33}\,\mathrm{kg^3})^{1/3}.
  \]
  \hfill(0.5)

  Thus,
  \[
    M_0^{\mathrm{max}} = 1.7 \times 10^{11}\,\mathrm{kg}.
  \]
  \hfill(0.5)

  Note that this was already found in (T12.2a), so the calculation above is
  not needed.

  \textbf{The smallest} $M_0$ will correspond to a BH that has evaporated as
  early as possible. Such a BH would have evaporated far away from Earth,
  and its radiation would be redshifted.

  The smallest time of evaporation would correspond to the largest redshift
  factor, and hence the largest $E_{\mathrm{emitted}}(t)$, which is
  $E_{\mathrm{emitted}}(t) = k_{\mathrm{B}} T_{\mathrm{Planck}}$.
  \hfill(1.0)

  The ratio of the scale factors at the time of the above PBH evaporation
  and now is
  \[
    \frac{a(t)}{a(t_{\mathrm{now}})}
    = \frac{E_{\mathrm{det}}}{k_{\mathrm{B}} T_{\mathrm{Planck}}}
    = \frac{3.0 \times 10^{20}\,\mathrm{eV}}{10^{19}\,\mathrm{GeV}}
    = 3.0 \times 10^{-8}
  \]
  With $a(t_{\mathrm{now}}) = 1.0$, we have
  $a(t) = 3.0 \times 10^{-8}$. \hfill(1.0)

  Using the plot of $a(t)$ vs $t$ from the figure, $t \simeq 10^{4}$\,s.
  Note that for this $t$, the second term in the original $M_0^3$ expression
  is still negligible. \hfill(1.0)

  Therefore,
  \[
    M_0^{\mathrm{min}} \approx (3At)^{1/3}
    \approx 4.9 \times 10^{6}\,\mathrm{kg}.
  \]
  \hfill(1.0)
\end{description}

\solhead{10.20}{Isolated black hole}
% source id: 2023_DA_Q2
\begin{description}
\item[(a)]
  \[
    \theta_{\mathrm{E}}
    = \sqrt{\frac{4GM}{c^2}\,\frac{D_s - D_\ell}{D_s D_\ell}}
    = \sqrt{\frac{4GM}{\mathrm{au}\,c^2}
      \left(\frac{\mathrm{au}}{D_\ell} - \frac{\mathrm{au}}{D_s}\right)}
    = 2.7\,\mathrm{mas}
  \]
  The angular resolution of modern large ($D \approx 10$\,m) optical
  ($\lambda = 550$\,nm) telescopes is
  $\theta_0 = 1.22\lambda/D \approx 14$\,mas. Thus,
  $\theta_0 \gg \theta_{\mathrm{E}}$, so the images created during
  microlensing events cannot be resolved by these telescopes.
  \hfill(2 points)

\item[(b)]
  \begin{align*}
    \theta_c &= \frac{\theta_+ A_+ + \theta_- A_-}{A_+ + A_-}
    = \frac{\frac{1}{4}\left(u + \sqrt{u^2+4}\right)
        \left(\frac{u^2+2}{u\sqrt{u^2+4}} + 1\right)
      + \frac{1}{4}\left(u - \sqrt{u^2+4}\right)
        \left(\frac{u^2+2}{u\sqrt{u^2+4}} - 1\right)}
      {\frac{u^2+2}{u\sqrt{u^2+4}}}\,\theta_{\mathrm{E}}\\
    &= \frac{\left(u + \sqrt{u^2+4}\right)
        \left(u^2 + 2 + u\sqrt{u^2+4}\right)
      + \left(u - \sqrt{u^2+4}\right)
        \left(u^2 + 2 - u\sqrt{u^2+4}\right)}
      {4(u^2+2)}\,\theta_{\mathrm{E}}\\
    &= \frac{2u(u^2+2) + 2u(u^2+4)}{4(u^2+2)}\,\theta_{\mathrm{E}}
    = \frac{2u(2u^2+6)}{4(u^2+2)}\,\theta_{\mathrm{E}}
    = \frac{u(u^2+3)}{u^2+2}\,\theta_{\mathrm{E}}
  \end{align*}
  \hfill(8 points)

\item[(c)]
  \[
    \Delta\theta = \theta_c - \theta
    = \frac{u(u^2+3)}{u^2+2}\,\theta_{\mathrm{E}} - u\theta_{\mathrm{E}}
    = \frac{u(u^2+3) - u(u^2+2)}{u^2+2}\,\theta_{\mathrm{E}}
    = \frac{u}{u^2+2}\,\theta_{\mathrm{E}}
  \]
  \hfill(3 points)
  \[
    \Delta\theta(u = 0) = 0
  \]
  so there is no deflection when the lens and the source are nearly
  perfectly aligned. \hfill(1 points)

\item[(d)] Plots of the measured E and N positions against time.
  \hfill(10 points, 5 points for each graph)
  \ioaafig{2023_DA_Q2-s1.pdf}

\item[(e)] I will use the fact that the first epoch of astrometric
  observations was taken close to the peak of the light curve (that is,
  $u_1 \approx 0$, that is, almost no astrometric deflection). Similarly,
  astrometric deflection is close to zero for the last epoch. Thus,
  \begin{align*}
    \mu_E &\approx \frac{x_{E,8} - x_{E,1}}{t_8 - t_1}
      = -2.247 \pm 0.029\,\mathrm{mas/yr}\\
    \mu_N &\approx \frac{x_{N,8} - x_{N,1}}{t_8 - t_1}
      = -3.674 \pm 0.038\,\mathrm{mas/yr}
  \end{align*}
  \hfill(8 points)

  \emph{No points should be given if a student does not recognize that the
  deflection is zero during the first epoch, e.g., they try fitting a
  straight line to all data points.}

\item[(f)] I fitted a straight line joining the first and last epoch data
  and then subtracted it from astrometric measurements. This is because the
  observed path of the source on the sky is the sum of two effects: the
  rectilinear proper motion of the source and the astrometric deflection:
  \begin{align*}
    x(E) &= x_1^E + (t - t_1)\mu_{\mathrm{E}} + \Delta\theta(E)\\
    x(N) &= x_1^N + (t - t_1)\mu_{\mathrm{N}} + \Delta\theta(N),
  \end{align*}
  where $x_1^E$ is the East position of the source during the first epoch,
  $x_1^N$ is the North position of the source during the first epoch, $t_1$
  is the time of the first observation, and $\Delta\theta(E)$ and
  $\Delta\theta(N)$ is the astrometric deflection due to microlensing in
  East and North direction, respectively.
  \ioaafig{2023_DA_Q2-s2.pdf}

  Thus, the astrometric deflection during $i$th epoch is:
  \begin{align*}
    \Delta\theta(E)_i &= x_i^E - x_1^E - (t_i - t_1)\mu_{\mathrm{E}}\\
    \Delta\theta(N)_i &= x_i^N - x_1^N - (t_i - t_1)\mu_{\mathrm{N}}
  \end{align*}
  and the total deflection is:
  \[
    \Delta\theta_i = \sqrt{\Delta\theta(E)_i^2 + \Delta\theta(N)_i^2}.
  \]
  I will also use the fact that $u_0 \approx 0$, so
  $u = (t - t_0)/t_E$, where $t_0 = 2455763$ is the peak time. Results of
  my calculations are shown in the table below.
  \ioaafig{2023_DA_Q2-s3.pdf}
  \ioaafig{2023_DA_Q2-s4.pdf}

  \begin{center}
  \begin{tabular}{rrrrr}
  \hline\hline
  Epoch & $u$ & $\Delta\theta$ (E,mas) & $\Delta\theta$ (N,mas)
    & $\Delta\theta$ (mas)\\
  \hline
  1 & 0.01 & $0.00 \pm 0.13$ & $-0.00 \pm 0.16$ & $0.00 \pm 0.21$\\
  2 & 0.42 & $0.36 \pm 0.12$ & $-0.84 \pm 0.24$ & $0.91 \pm 0.27$\\
  3 & 1.69 & $0.43 \pm 0.14$ & $-1.50 \pm 0.08$ & $1.56 \pm 0.16$\\
  4 & 1.75 & $0.95 \pm 0.36$ & $-1.40 \pm 0.77$ & $1.69 \pm 0.85$\\
  5 & 2.69 & $0.47 \pm 0.21$ & $-1.58 \pm 0.12$ & $1.65 \pm 0.24$\\
  6 & 3.34 & $0.44 \pm 0.07$ & $-0.90 \pm 0.40$ & $1.00 \pm 0.41$\\
  7 & 4.83 & $0.21 \pm 0.25$ & $-0.47 \pm 0.29$ & $0.51 \pm 0.38$\\
  8 & 9.04 & $0.00 \pm 0.12$ & $-0.00 \pm 0.17$ & $0.00 \pm 0.21$\\
  \hline
  \end{tabular}
  \end{center}
  \hfill(20 points)

\item[(g)] We would like to fit the function
  $\Delta\theta = \frac{u}{u^2+2}\theta_{\mathrm{E}}$ to the data. Thus, my
  ``new'' independent variable would be $x' = u/(u^2+2)$. Now, I would like
  to fit the function $y' = \theta_{\mathrm{E}} x'$, where
  $y' = \Delta\theta$. Thus
  \[
    \theta_{\mathrm{E}}
    = \frac{\sum y'_i x'_i/\sigma_i^2}{\sum x_i^{\prime 2}/\sigma_i^2}
      \pm \frac{1}{\sqrt{\sum x_i^{\prime 2}/\sigma_i^2}}
    = 4.5 \pm 0.4\,\mathrm{mas}
  \]
  \ioaafig{2023_DA_Q2-s5.pdf}

  Alternative solution: The gravitational deflection is described by the
  formula $\Delta\theta = \frac{u}{u^2+2}\theta_{\mathrm{E}}$. It can be
  demonstrated that this function reaches a local maximum for
  $u = \sqrt{2}$ with
  $\Delta\theta_{\mathrm{max}} = \frac{\sqrt{2}}{4}\theta_{\mathrm{E}}
  = 0.354\,\theta_{\mathrm{E}}$. From the data, we can estimate the maximum
  deflection of $\Delta\theta_{\mathrm{max}} = 1.59 \pm 0.14$\,mas. Thus,
  $\theta_{\mathrm{E}} = 4.5 \pm 0.4$\,mas. \hfill(16 points)

\item[(h)] Let $\pi_{\mathrm{rel}} = \pi_l - \pi_s$ be the relative
  lens--source parallax. From the definition of the angular Einstein radius
  we have
  \[
    \theta_{\mathrm{E}}
    = \sqrt{\frac{4GM}{c^2}\,\frac{D_s - D_l}{D_s D_l}}
    = \sqrt{\frac{4GM}{c^2}\left(\frac{1}{D_l} - \frac{1}{D_s}\right)}
    = \sqrt{\frac{4GM}{c^2\,\mathrm{au}}
      \left(\frac{\mathrm{au}}{D_l} - \frac{\mathrm{au}}{D_s}\right)}
    = \sqrt{\frac{4GM\pi_{\mathrm{rel}}}{c^2\,\mathrm{au}}}.
  \]
  From the definition of the microlensing parallax
  $\pi_{\mathrm{rel}} = \pi_{\mathrm{E}}\theta_E$, so this equation
  becomes:
  \[
    \theta_{\mathrm{E}}
    = \sqrt{\frac{4GM\pi_{\mathrm{E}}\theta_{\mathrm{E}}}{c^2\,\mathrm{au}}},
  \]
  and hence
  \[
    M = \frac{\theta_{\mathrm{E}} c^2\,AU}{4G\pi_{\mathrm{E}}}
    = 5.8 \pm 0.8\,M_\odot.
  \]
  The uncerainty% sic
  {} on $M$ can be determined from the relation:
  \[
    \frac{\Delta M}{M}
    = \sqrt{\left(\frac{\Delta\theta_E}{\theta_{\mathrm{E}}}\right)^2
      + \left(\frac{\Delta\pi_{\mathrm{E}}}{\pi_{\mathrm{E}}}\right)^2}
  \]
  \hfill(7 points)
\end{description}

